\section{Theme Codebook}
\label{app:codebook}
\label{app:codebook-table}
This section reproduces the consolidated codebook from Step~2 of the formative analysis (Chapter~\ref{chap:formative-study}, Section~\ref{sec:theme-consolidation}): the 64 themes, T1 to T64, into which the 435 initial codes were grouped, and beneath each theme the excerpts that informed it, 437 in all.

A theme's header row carries its identifier, its need label, $n$, and its disposition. Here $n$ counts the interviews in which the theme recurred, not the excerpts beneath it, which is usually the larger number because one participant could speak to a theme several times. The disposition is the verdict of the scope review (Section~\ref{sec:scope-refinement}): carried forward into the requirements derivation of Section~\ref{chap:from-themes-to-requirements}, excluded together with the boundary that excluded it, or methodological artefact for a theme about the interview situation rather than a planning need. Carried forward records that verdict only, not that the requirements built on the theme were implemented.

Each row beneath a header is one excerpt: the interview it came from (I1 to I6), the author's English translation, and the German original, so that every translation can be checked against what was said.

% Generated by tools/build_codebook.py -- do not edit by hand.
\small
\begin{longtable}{@{}p{0.7cm}p{6.0cm}p{6.0cm}@{}}
\toprule
\textbf{Int.} & \textbf{Excerpt (translated)} & \textbf{Excerpt (German original)} \\
\midrule
\endfirsthead
\toprule
\textbf{Int.} & \textbf{Excerpt (translated)} & \textbf{Excerpt (German original)} \\
\midrule
\endhead
\bottomrule
\endfoot
\multicolumn{3}{@{}p{12.9cm}@{}}{\textbf{T1} --- Needs multi-level study planning overview (semester → future semesters → whole programme) with visual progress \quad ($n=5$; Carried forward)} \\[2pt]
I1 & \enquote{…three levels… for my current semester, my deadlines and courses… then… next semester or the next few semesters. And then the third level… how am I on track for my entire degree, in percent or visually…} & \enquote{…drei Ebenen… für mein aktuelles Semester mir meine Deadlines und Kurse… dann… nächstes Semester oder die nächste paar Semester. Und dann die dritte Ebene… wie bin ich on track für mein komplettes Studium in Prozent oder visuell…} \\
I2 & \enquote{The whole degree… through to graduation now, to the title… to the Bachelor's title or to the Master's.} & \enquote{Das ganze Studium… bis zum Abschluss jetzt, bis zum Titel… bis zum Bachelortitel oder bis zum Master.} \\
I4 & \enquote{Semester view. / plans their Master's, their Bachelor's.} & \enquote{Semestersicht.} / \enquote{sein Master, seinen Bachelor plant.} \\
I5 & \enquote{the whole degree is planned… which courses do I take in which semester, and why?} & \enquote{das ganze Studium geplant wird… welche Kurse nehme ich in welchem Semester und warum?} \\
I5 & \enquote{only for my current semester.} & \enquote{nur für mein jetziges Semester.} \\
I5 & \enquote{plan the whole Master's.} & \enquote{den ganzen Master plane.} \\
I6 & \enquote{how you… divide the curriculum across the semesters.} & \enquote{wie man sich… das Curriculum auf den Semestern aufteilt.} \\
\addlinespace
\multicolumn{3}{@{}p{12.9cm}@{}}{\textbf{T2} --- Needs granular course-level visibility (appointments, deadlines, exam dates) \quad ($n=3$; Excluded, OOS2)} \\[2pt]
I1 & \enquote{And maybe even… really granular… per course, so that I know where my appointments and my deadlines are…} & \enquote{Und vielleicht sogar… ganz Granulat… pro Lehrveranstaltung, dass ich da weiß, wo sind da meine Termine und meine Deadlines…} \\
I4 & \enquote{just the appointments.} & \enquote{die Termine halt.} \\
I5 & \enquote{Calendar notification… two weeks beforehand.} & \enquote{Kalender Notification… zwei Wochen davor.} \\
\addlinespace
\multicolumn{3}{@{}p{12.9cm}@{}}{\textbf{T3} --- Needs clear scope/definition of what planning includes (deadlines vs course list) \quad ($n=2$; Methodological artefact, clarifying the scope of the interview)} \\[2pt]
I1 & \enquote{But does that mean with deadlines, with exam dates, or simply: in this semester I've five courses?} & \enquote{Das bedeutet aber dann mit Deadlines, mit Testterminen oder einfach nur in dem Semester habe ich fünf Kurse?} \\
I6 & \enquote{we are focusing only on courses… when they are chosen.} & \enquote{wir konzentrieren uns nur auf Kurse… wann diese gewählt werden.} \\
\addlinespace
\multicolumn{3}{@{}p{12.9cm}@{}}{\textbf{T4} --- Needs planning that accounts for personal constraints and supports rescheduling/adjustment \quad ($n=5$; Carried forward)} \\[2pt]
I1 & \enquote{Because I always worked alongside, it never worked out. That means I had to start shifting things a little…} & \enquote{Weil ich immer nebenbei gearbeitet habe, ging das nie auf. Das heißt, ich musste dann immer schon ein bisschen anfangen zu verschieben…} \\
I1 & \enquote{Because I always worked alongside… had to… shift things… but I just got on with it.} & \enquote{Weil ich immer nebenbei gearbeitet habe… musste… verschieben… aber habe es einfach gemacht.} \\
I3 & \enquote{because there are some subjects I didn't manage in the second semester.} & \enquote{weil es einige Fächer gibt, die ich doch nicht geschafft habe im zweiten Semester.} \\
I4 & \enquote{but that's a part-time job.} & \enquote{das ist aber ein Teilzeitjob.} \\
I5 & \enquote{I'd rather… manage 100\% and still be able to live a little.} & \enquote{lieber… hundertprozentig schaffe und auch noch ein bisschen leben kann.} \\
I5 & \enquote{maybe I'll try to shift some things around.} & \enquote{ich versuch vielleicht was umzuschichten.} \\
I6 & \enquote{what do I have to replan.} & \enquote{was muss ich umplanen.} \\
\addlinespace
\multicolumn{3}{@{}p{12.9cm}@{}}{\textbf{T5} --- Needs structure/guidance when programme shifts from fixed curriculum to free choice \quad ($n=5$; Carried forward)} \\[2pt]
I1 & \enquote{The first two, three semesters… I simply lived by the plan, or tried to stick to it.} & \enquote{Die ersten zwei, drei Semester… habe ich nach dem Plan einfach gelebt oder den versucht einzuhalten.} \\
I1 & \enquote{In the Master's… only the first semester, and after that free choice.} & \enquote{Im Master… nur das erste Semester und danach freie Wahl.} \\
I2 & \enquote{The first semester was simply what was prescribed… the standard package with 29~ECTS.} & \enquote{Das erste Semester war… einfach das, was vorgegeben wurde… das Standardpaket mit neunundzwanzig ECTS.} \\
I2 & \enquote{Second semester… almost identical… one module was optional.} & \enquote{Zweites Semester… ziemlich identisch… ein Modul war freiwillig.} \\
I2 & \enquote{From the third onwards I really did draw up my timetable… semester plan.} & \enquote{Ab dem dritten habe ich dann wirklich meinen Stundenplan… Semesterplan drüber erstellt.} \\
I3 & \enquote{in the first year I only took the subjects that… I had to do in the plan.} & \enquote{im ersten Jahr habe ich eigentlich nur die Fächer genommen, die… im Plan machen musste.} \\
I5 & \enquote{for the first semester… just as it says in the curriculum.} & \enquote{für das erste Semester… so wie es im Studienplan steht.} \\
I6 & \enquote{I'm studying in the Master's… Software Engineering.} & \enquote{Ich studiere im Master… Software Engineering.} \\
\addlinespace
\multicolumn{3}{@{}p{12.9cm}@{}}{\textbf{T6} --- Needs support for course choice + feasibility planning (content-/interest-driven, what's doable) \quad ($n=6$; Carried forward)} \\[2pt]
I1 & \enquote{At the beginning I simply chose the courses that appealed to me… I can manage that this semester, but without much advance planning.} & \enquote{Am Anfang habe ich einfach die Kurse gewählt, die mich angesprochen haben… das schaffe ich das Semester, aber ohne große Vorplanung.} \\
I2 & \enquote{In terms of content… it works itself out.} & \enquote{Inhaltlich… ergibt sich selber.} \\
I3 & \enquote{you definitely go at it by content.} & \enquote{du gehst inhaltlich auf jeden Fall ran.} \\
I3 & \enquote{otherwise I tend to take subjects that interest me.} & \enquote{sonst nehme ich eigentlich eher Fächer, die mich interessieren.} \\
I4 & \enquote{which courses interest me.} & \enquote{welche Lehrveranstaltungen interessieren mich.} \\
I4 & \enquote{I'm specialising in Security.} & \enquote{ich mach Security schwerpunktmäßig.} \\
I4 & \enquote{Software Engineering… not that important.} & \enquote{Software Engineering… nicht so wichtig.} \\
I5 & \enquote{Computer graphics… took it because it was just cool.} & \enquote{Computergrafik… genommen, weil es halt cool war.} \\
I6 & \enquote{so that you don't do anything unnecessary.} & \enquote{damit man eben nichts macht, was unnötig ist.} \\
I6 & \enquote{see at a glance… what is actually possible.} & \enquote{auf einen Blick sehen… was ist überhaupt möglich.} \\
I6 & \enquote{doesn't know whether they will be credited.} & \enquote{nicht weiß, ob man die angerechnet bekommt.} \\
\addlinespace
\multicolumn{3}{@{}p{12.9cm}@{}}{\textbf{T7} --- Needs a triggered \enquote{bring structure} moment to be supported (timing of when planning becomes important) \quad ($n=5$; Carried forward)} \\[2pt]
I1 & \enquote{From around the third semester you realise: `Okay, now I need to start bringing some structure in.'} & \enquote{So ab dem dritten Semester merkt man irgendwann: `Okay, jetzt muss ich langsam mal irgendwie Struktur reinbringen.'} \\
I3 & \enquote{at the moment it isn't that severe… not really a problem.} & \enquote{da ist jetzt im Moment noch nicht so heftig viel… wirklich ein Problem.} \\
I4 & \enquote{at the end of the degree.} & \enquote{am Ende des Studiums.} \\
I5 & \enquote{in the holidays… enter it for the third semester.} & \enquote{in den Ferien… für das dritte Semester eintragen.} \\
I5 & \enquote{only did it much later, unfortunately.} & \enquote{erst viel später gemacht, leider.} \\
I5 & \enquote{it would have been better if I'd done it earlier.} & \enquote{das wäre besser, wenn ich es früher gemacht hätte.} \\
I5 & \enquote{instead of… from semester to semester.} & \enquote{anstatt… von Semester zu Semester.} \\
I6 & \enquote{start and end of the semester.} & \enquote{Beginn und Ende des Semesters.} \\
\addlinespace
\multicolumn{3}{@{}p{12.9cm}@{}}{\textbf{T8} --- Needs features that reduce overwhelm and emotional burden from planning \quad ($n=2$; Carried forward)} \\[2pt]
I1 & \enquote{And that was also very depressing, because only then did I see how much I still lack.} & \enquote{Und das war auch sehr deprimierend, weil da habe ich dann erst gesehen, wie viel mir eigentlich noch fehlt.} \\
I1 & \enquote{At the start… short-term… at some point I panicked… drew up a plan that became more and more concrete…} & \enquote{Am Anfang… kurzfristig… irgendwann habe ich eine Panik geschoben… Plan aufgestellt, der immer konkreter wurde…} \\
I5 & \enquote{I was completely worn out.} & \enquote{ich war so fix und fertig.} \\
\addlinespace
\multicolumn{3}{@{}p{12.9cm}@{}}{\textbf{T9} --- Needs institutional overview that is actually usable (uni system currently clunky / not computing / not clear) \quad ($n=5$; Carried forward)} \\[2pt]
I1 & \enquote{…in TISS… created categories, fifth semester, sixth semester… tried to sort my courses… into them.} & \enquote{...in TISS… Kategorien angelegt, so fünfte Semester, sechste Semester… versucht, meine Kurse… reinzuordnen.} \\
I2 & \enquote{the overview isn't clear… it doesn't calculate anything.} & \enquote{die Übersicht ist… nicht übersichtlich… tut nichts berechnen.} \\
I3 & \enquote{sometimes you can't remove something any more… not really meaningful.} & \enquote{manchmal kann man was nicht mehr entfernen… nicht so richtig aussagekräftig.} \\
I3 & \enquote{curriculum in TISS… not that good.} & \enquote{Studienplan in TISS… nicht so gut.} \\
I4 & \enquote{I don't have to log into TISS.} & \enquote{muss ich mich nicht im TISS einloggen.} \\
I6 & \enquote{outdated UI… not trivial.} & \enquote{veraltete UI… nicht trivial.} \\
\addlinespace
\multicolumn{3}{@{}p{12.9cm}@{}}{\textbf{T10} --- Needs to work with / unify tools \& sources (uni system + lightweight personal methods) \quad ($n=6$; Carried forward)} \\[2pt]
I1 & \enquote{…tools… Excel… Notes… TISS… more or less an institutional tool…} & \enquote{...Tools… Excel… Notes… TIS[S]… mehr oder weniger ein institutionelles Tool…} \\
I2 & \enquote{the good old pen and paper… and the Excel sheet.} & \enquote{der gute, das alte Stift, Papier… und halt die Excel-Tabelle.} \\
I3 & \enquote{I look at TISS… what is on offer in the winter semester.} & \enquote{ich guck auf TIS… was es eben im Wintersemester gibt.} \\
I3 & \enquote{you can also create these favourites… I use that.} & \enquote{diese Favoriten kann man auch erstellen… das verwende ich.} \\
I3 & \enquote{No, not really.} & \enquote{Nein, eigentlich nicht wirklich.} \\
I4 & \enquote{what really helped me was an Excel document.} & \enquote{was mir eigentlich sehr geholfen hat, war Excel-Dokument.} \\
I5 & \enquote{Studo? Does that count?} & \enquote{Studo? Zählt das dazu?} \\
I5 & \enquote{on TISS… under favourites.} & \enquote{auf TIS… unter Favoriten.} \\
I6 & \enquote{TISS for registering… favouriting.} & \enquote{TIS zum Anmelden… Favorisieren.} \\
\addlinespace
\multicolumn{3}{@{}p{12.9cm}@{}}{\textbf{T11} --- Needs peer-validated planning artefacts and reusable templates \quad ($n=1$; Carried forward)} \\[2pt]
I1 & \enquote{…a friend… sent me an Excel file that he had planned with…} & \enquote{...ein Freund… einen Excel-File geschickt, mit dem er… geplant hat.} \\
I1 & \enquote{…no finished Excel from the university, unfortunately… you have to manage it yourself somehow.} & \enquote{...kein fertiges Excel… von der Uni leider… sondern du musst das irgendwie selber managen.} \\
\addlinespace
\multicolumn{3}{@{}p{12.9cm}@{}}{\textbf{T12} --- Needs a planning view that improves on the institutional UI (visualisation + notes + grades) \quad ($n=2$; Carried forward)} \\[2pt]
I1 & \enquote{…just… a nicer view than in TISS, where you can also enter your grades and the hours, and where you can visualise things…} & \enquote{...nur… eine schönere Ansicht als in TISS, wo man auch die Noten eintragen kann und die Stunden und man sich auch visualisieren kann…} \\
I4 & \enquote{Bars… chart… shown with colours.} & \enquote{Säulen… Diagramm… mit Farben dargestellt.} \\
\addlinespace
\multicolumn{3}{@{}p{12.9cm}@{}}{\textbf{T13} --- Needs system-aware requirement checking (rules/requirements encoded) instead of manual tracking \quad ($n=6$; Carried forward)} \\[2pt]
I1 & \enquote{…went through the requirements again and again… hard requirements… entered them into the Excel…} & \enquote{...immer wieder auch die Requirements durchgegangen… hard Requirements… in das Excel eingetragen…} \\
I1 & \enquote{…Excel doesn't know your requirements… it doesn't detect errors either.} & \enquote{...excel halt nicht deine Requirements kennt… erkennt auch keine Fehler.} \\
I1 & \enquote{…you have to double and triple check everything and hope it fits…} & \enquote{...musst alles doppelt, dreifach überprüfen, hoffen, dass es passt…} \\
I2 & \enquote{keeping the rules in your head, what actually has to be done.} & \enquote{die Regeln im Kopf zu behalten, was überhaupt sein muss.} \\
I2 & \enquote{it doesn't show me what I'm missing.} & \enquote{es zeigt mir nicht, was mir fehlt.} \\
I2 & \enquote{seven out of ten narrow-elective subjects… required to do.} & \enquote{sieben von zehn Wahlpflichtfächern… verpflichtet zu machen.} \\
I2 & \enquote{open the curriculum… written so terribly.} & \enquote{Curriculum öffnen… so schrecklich geschrieben.} \\
I2 & \enquote{before you… look at modules… you have to look at, okay, what do you have to do?} & \enquote{bevor man… Module anschaut… muss man sich anschauen, okay, was muss man machen?} \\
I3 & \enquote{narrow elective… where you have to do at least seven subjects.} & \enquote{enge Wahl… wo man zumindest sieben Fächer machen muss.} \\
I4 & \enquote{which subjects do I need… which compulsory subjects.} & \enquote{welche Fächer brauch ich… welche Pflichtfächer.} \\
I4 & \enquote{to see: what do you have to do?} & \enquote{zu sehen: Was musst du machen?} \\
I5 & \enquote{I had to write to the dean's office.} & \enquote{musste ich das Dekanat anschreiben.} \\
I5 & \enquote{an Excel sheet… really, really good.} & \enquote{ein Excel-Sheet… richtig, richtig gut.} \\
I5 & \enquote{not sure any more… whether I should take something.} & \enquote{nicht mehr sicher… ob ich was nehmen soll.} \\
I5 & \enquote{check whether you meet all the requirements.} & \enquote{checken, ob man quasi alle Requirements erfüllt.} \\
I5 & \enquote{read this huge PDF.} & \enquote{dieses große PDF… lesen.} \\
I6 & \enquote{credited as interdisciplinary.} & \enquote{angerechnet werden als interdisziplinäre.} \\
I6 & \enquote{doesn't know whether they will be credited.} & \enquote{nicht weiß, ob man die angerechnet bekommt.} \\
I6 & \enquote{interdisciplinary courses… maximum number.} & \enquote{interdisziplinäre Kurse… maximale Anzahl.} \\
\addlinespace
\multicolumn{3}{@{}p{12.9cm}@{}}{\textbf{T14} --- Needs automatic updates / low-maintenance upkeep when curriculum changes \quad ($n=2$; Excluded, OOS3)} \\[2pt]
I1 & \enquote{…because it was outdated and… the curriculum changed, I kept going through the requirements again…} & \enquote{...weil das aber outdatet war und… Studienplan geändert hat, bin ich dann immer wieder… Requirements durchgegangen…} \\
I1 & \enquote{Because I don't trust myself… in six months… `what did I do there'… I have to think it all through again.} & \enquote{Weil ich traue mir selber nicht… in einem halben Jahr… `was habe ich da gemacht'… muss alles wieder neu bedenken.} \\
I5 & \enquote{the transitional provisions.} & \enquote{die Übergangsbestimmungen.} \\
I5 & \enquote{very, very large.} & \enquote{sehr, sehr groß.} \\
\addlinespace
\multicolumn{3}{@{}p{12.9cm}@{}}{\textbf{T15} --- Needs integrated tool aligned with university requirements + personal needs \quad ($n=2$; Carried forward)} \\[2pt]
I1 & \enquote{…the lack of a tool that… is integrated with what the university needs… and what you personally… need…} & \enquote{...Fehlen von einem Tool, das… integriert ist in dem, was die Uni braucht… und was auch du persönlich… brauchst…} \\
I4 & \enquote{a plan… individual for everyone… effective… realistic… to carry out.} & \enquote{ein Plan… für jeden individuell… effektiv… realitätsnah… umzusetzen.} \\
\addlinespace
\multicolumn{3}{@{}p{12.9cm}@{}}{\textbf{T16} --- Needs progress tracking + status dashboard (credits, \%, missing modules, end date visibility) \quad ($n=5$; Carried forward)} \\[2pt]
I1 & \enquote{…through that I saw: how much do I still need?… I'm still missing ten… I still need…} & \enquote{...dadurch gesehen: Wie viel brauche ich noch?… mir fehlen noch zehn… brauche ich noch…} \\
I1 & \enquote{A year ago I knew… that my degree would end in February… Two weeks before that, I had no idea at all.} & \enquote{Vor einem Jahr wusste ich… dass jetzt mit Februar mein Studium enden wird… Also zwei Wochen davor noch gar nicht.} \\
I1 & \enquote{…I still don't see any status, where am I right now?} & \enquote{...ich sehe immer noch keinen Status, wo bin ich gerade?} \\
I1 & \enquote{…how many percent… am I still missing?… now I've sixty percent…} & \enquote{...wie viel Prozent… fehlen mir noch?… jetzt habe ich sechzig Prozent…} \\
I2 & \enquote{it would be… nice if there were an overview like that for the Bachelor's.} & \enquote{es wäre… schön, wenn es so eine Übersicht gäbe für den Bachelor.} \\
I2 & \enquote{an overview… progress bar… showing that I've already done those.} & \enquote{eine Übersicht… Progress-Bahn… dass ich schon die gemacht habe.} \\
I2 & \enquote{scrolling through all my transcripts.} & \enquote{durch meine ganzen Zeugnisse durchscrollen.} \\
I2 & \enquote{where do you stand now… what is actively missing?} & \enquote{wo steht man jetzt… was fehlt denn aktiv?} \\
I2 & \enquote{three of seven have been done.} & \enquote{drei von sieben wurden gemacht.} \\
I4 & \enquote{what have I completed… what ECTS am I still missing.} & \enquote{was hab ich abgeschlossen… was fehlt mir noch für ECTS.} \\
I4 & \enquote{per module a bar.} & \enquote{je Modul halt eine Säule.} \\
I4 & \enquote{the thirty for the diploma thesis and three from the seminar.} & \enquote{die dreißig Diplomarbeit und drei vom Seminar.} \\
I5 & \enquote{twelve of eighteen ECTS done.} & \enquote{zwölf von achtzehn ECTS fertig.} \\
I6 & \enquote{always an up-to-date status.} & \enquote{immer einen aktuellen Status.} \\
I6 & \enquote{how much have I completed so far.} & \enquote{wie viel habe ich aktuell abgeschlossen.} \\
I6 & \enquote{overview… of my ECTS.} & \enquote{Überblick… über meine ECTS.} \\
I6 & \enquote{you don't lose track so easily.} & \enquote{verliert nicht so leicht den Überblick.} \\
I6 & \enquote{see at a glance.} & \enquote{auf einen Blick sehe.} \\
\addlinespace
\multicolumn{3}{@{}p{12.9cm}@{}}{\textbf{T17} --- Needs automatic requirement-fulfilment feedback (overall + categories) \quad ($n=1$; Carried forward)} \\[2pt]
I1 & \enquote{…that tells me… you've already fulfilled five of seven requirements… all transferable skills done.} & \enquote{...dass mir sagt… hast du jetzt schon fünf von sieben Requirements erfüllt… alle Transfrarable Skills erledigt.} \\
I1 & \enquote{…`I've already fulfilled 20\% of the core modules…'} & \enquote{...`ich habe jetzt schon 20\% der Core Modelle erfüllt…'} \\
\addlinespace
\multicolumn{3}{@{}p{12.9cm}@{}}{\textbf{T18} --- Needs future-state preview + completion forecasting \quad ($n=2$; Carried forward)} \\[2pt]
I1 & \enquote{…if I add that, then in the future I would have 25\% fulfilled.} & \enquote{...wenn ich das… zufüge, dann in Zukunft hätte ich 25\% erfüllt.} \\
I1 & \enquote{…if I do that in this semester… then I've… reached 100\% in the seventh semester…} & \enquote{...wenn ich in dem Semester das mache… dann habe ich… im siebten Semester 100\% erreicht…} \\
I4 & \enquote{how quickly you want to finish.} & \enquote{wie schnell man fertig werden will.} \\
\addlinespace
\multicolumn{3}{@{}p{12.9cm}@{}}{\textbf{T19} --- Needs a clear distinction planned vs completed, with intuitive completion visuals + motivating feedback \quad ($n=5$; Carried forward)} \\[2pt]
I1 & \enquote{What have I chosen?… what have I already completed? Separately.} & \enquote{Was habe ich gewählt?… was habe ich schon abgeschlossen? Separat.} \\
I1 & \enquote{A star means finished, for me.} & \enquote{Stern heißt für mich… fertig.} \\
I1 & \enquote{Maybe grey them out as well… you no longer need to highlight it visually.} & \enquote{Vielleicht auch ausgrauen… du brauchst es nicht mehr visuell hervorheben.} \\
I1 & \enquote{…satisfying when it becomes less and less.} & \enquote{...satisfying, wenn das immer weniger wird.} \\
I1 & \enquote{That is another reward… I need a reward.} & \enquote{Das ist dann noch ein Reward… ich brauche noch einen Reward.} \\
I2 & \enquote{Differentiation… planned… completed.} & \enquote{Differenzierung… geplant… abgeschlossen.} \\
I2 & \enquote{once it's passed… I don't really need to see it any more.} & \enquote{wenn es bestanden ist… brauche ich es eigentlich nicht mehr zu sehen.} \\
I2 & \enquote{Plus and star… easier to understand.} & \enquote{Plus-Sternchen… verständlicher.} \\
I4 & \enquote{coloured in… which ones am I attending?} & \enquote{eingefärbt… welche besuche ich da?} \\
I4 & \enquote{Distinction between… planned… already done.} & \enquote{Unterscheidung zwischen… geplant… schon gemacht.} \\
I5 & \enquote{can see whether it's already completed.} & \enquote{sehen kann, ob es schon abgeschlossen ist.} \\
I5 & \enquote{coloured… done.} & \enquote{gefärbt… done.} \\
I5 & \enquote{without them… disappearing altogether.} & \enquote{ohne dass sie… nicht mehr auftauchen.} \\
I5 & \enquote{added… finished.} & \enquote{hinzugefügt… fertig.} \\
I6 & \enquote{completed… in progress.} & \enquote{abgeschlossen… in progress.} \\
I6 & \enquote{that's how I plan in my Excel sheet.} & \enquote{so plane ich halt in meinem Excel-Sheet.} \\
I6 & \enquote{different stages.} & \enquote{verschiedene Stadi.} \\
I6 & \enquote{marked by colour… done.} & \enquote{farblich markiert… erledigt.} \\
I6 & \enquote{see at a glance.} & \enquote{auf einen Blick sehe.} \\
\addlinespace
\multicolumn{3}{@{}p{12.9cm}@{}}{\textbf{T20} --- Needs an exploration-first workflow (show all → select → then plan semesters) \quad ($n=5$; Carried forward)} \\[2pt]
I1 & \enquote{…cool as a first step… not thinking yet about which semester… but rather…} & \enquote{...cool als erster Schritt… noch gar keine Gedanken, welches Semester… sondern…} \\
I1 & \enquote{…final tree, and then I do the semester view.} & \enquote{...finalen Baum und dann mache ich die Semesteransicht.} \\
I1 & \enquote{…Show all… an immediate picture… satisfying… if… overwhelmed… then at least I know that this is all of it.} & \enquote{...Show all… sofort ein Bild… satisfying… wenn… overwhelmed… dann weiß ich, das ist alles.} \\
I1 & \enquote{…I click plus, plus, plus… and then I've all the ones that interest me…} & \enquote{...sage plus, plus, plus… und dann habe ich einmal alle, die mich interessieren…} \\
I1 & \enquote{…back into… the selected ones… look… at how that feels…} & \enquote{...zurück in… Ausgewählten… schaue… wie sich das so anfühlt…} \\
I1 & \enquote{…less FOMO that I missed something again…} & \enquote{...weniger FOMO, dass ich wieder was verpasse…} \\
I2 & \enquote{at the beginning of the degree… then you see… all the subjects.} & \enquote{am Anfang des Studiums… dann sieht man… die ganzen Fächer.} \\
I4 & \enquote{Graph used for exploration, not planning / find out… what is there?} & \enquote{Graph used for exploration, not planning} / \enquote{rausfinden… was gibt's?} \\
I4 & \enquote{Discovery view perceived as valuable / That is cool.} & \enquote{Discovery view perceived as valuable} / \enquote{Das ist cool.} \\
I5 & \enquote{That is cool.} & \enquote{Das ist cool.} \\
I6 & \enquote{it simply shows me everything as a graph.} & \enquote{stellt mir einfach alles als Graphen dar.} \\
I6 & \enquote{for an overview… practical.} & \enquote{für einen Überblick… praktisch.} \\
I6 & \enquote{you don't lose track so easily.} & \enquote{verliert nicht so leicht den Überblick.} \\
I6 & \enquote{select in the graph… push it over into the plan.} & \enquote{im Graphen auswählen… in den Plan rüberschiebt.} \\
I6 & \enquote{pull it out… into a sub-space.} & \enquote{in einem Sub-Space… rausziehen.} \\
\addlinespace
\multicolumn{3}{@{}p{12.9cm}@{}}{\textbf{T21} --- Needs course pool + drag/drop planning while preserving spatial layout \quad ($n=4$; Carried forward)} \\[2pt]
I1 & \enquote{…pool… with the courses… where I can move them around…} & \enquote{...Pool… mit den Lehrveranstaltungen… wo ich das verschieben kann…} \\
I2 & \enquote{the position… stays.} & \enquote{die Position… bleibt.} \\
I2 & \enquote{like a desk… you just push things around.} & \enquote{wie ein Schreibtisch… schiebt halt die Sachen.} \\
I4 & \enquote{you track these items across.} & \enquote{man trackt diese Items rüber.} \\
I4 & \enquote{Drag and drop somehow.} & \enquote{Drag and Drop irgendwie.} \\
I6 & \enquote{select in the graph… push it over into the plan.} & \enquote{im Graphen auswählen… in den Plan rüberschiebt.} \\
\addlinespace
\multicolumn{3}{@{}p{12.9cm}@{}}{\textbf{T22} --- Needs automatic plan generation after selection (less manual dragging), respecting constraints and typical paths \quad ($n=1$; Excluded, OOS1)} \\[2pt]
I1 & \enquote{…Generate Plan, and then it fills them into the semesters for me automatically…} & \enquote{...Generate Plan und dann füllt's mir automatisch in die Semester…} \\
I1 & \enquote{…so that… the quota is met… how many hours I've.} & \enquote{...damit… Kontingent erfüllt wird… wie vielen Stunden ich habe.} \\
I1 & \enquote{…suggested… what do other students typically do… what does the university suggest.} & \enquote{...vorgeschlagen… was machen andere Studierende typischerweise… was wird von der Uni vorgeschlagen.} \\
I1 & \enquote{…then I don't have to drag anything around any more.} & \enquote{...dann brauche ich nichts mehr rumzuziehen.} \\
\addlinespace
\multicolumn{3}{@{}p{12.9cm}@{}}{\textbf{T23} --- Needs recommendations that are trustworthy and user-relevant (not generic; grounded; can use peer data carefully) \quad ($n=6$; Carried forward)} \\[2pt]
I1 & \enquote{Recommendations… I'm always very critical of recommendations, because the system can't know what my life is like… how many hours I work…} & \enquote{Empfehlungen… immer sehr kritisch, weil das System halt nicht wissen kann, wie mein Leben ist… wie viele Stunden ich arbeite…} \\
I1 & \enquote{…professors… describe courses well… what it says in there… often not what you… get.} & \enquote{...Profs… Lehrveranstaltungen gut beschreiben… was drin steht… nicht oft das, was man… bekommt.} \\
I1 & \enquote{…good student feedback… would have to feed into it.} & \enquote{...müsste… gutes Studierenden-Feedback… einfließen.} \\
I1 & \enquote{…very living… a living, well-documented Wikipedia… then I can trust it…} & \enquote{...very living… gut dokumentiertes Wikipedia… dann kann ich dem vertrauen…} \\
I1 & \enquote{…I could manage that myself even without a tool… `Show All… complete picture, and then I just work my way through it.'} & \enquote{...das würde ich auch ohne Tool selber schaffen… `Show All… komplettes Bild und dann ackere ich mich halt durch.'} \\
I1 & \enquote{Really good. Then I've the advantage again that I move through with other students.} & \enquote{Voll gut. Dann habe ich wieder den Vorteil, dass ich dann auch mit Studierenden mitgehe.} \\
I2 & \enquote{not really… recommendations that are based on nothing.} & \enquote{nicht wirklich… Empfehlungen, die auf nichts basieren.} \\
I2 & \enquote{content-based… dependency-based.} & \enquote{inhaltsbasiert… abhängigkeitsbasiert.} \\
I2 & \enquote{tags/interests… filtered out.} & \enquote{Tags/Interessen… rausgefiltert.} \\
I2 & \enquote{I can imagine that working well.} & \enquote{kann mir vorstellen, dass das gut funktioniert.} \\
I2 & \enquote{I'm the first with my curriculum.} & \enquote{ich bin der Erste mit meinem Curriculum.} \\
I2 & \enquote{they can use it from me.} & \enquote{die können von mir nutzen.} \\
I3 & \enquote{Nothing really comes to mind.} & \enquote{Mir fällt eigentlich nicht wirklich was ein.} \\
I3 & \enquote{I want what interests me.} & \enquote{ich will ja das, was mich interessiert.} \\
I3 & \enquote{help… with what I take first and what later.} & \enquote{helfen… was ich zuerst nehme und was später.} \\
I3 & \enquote{people… who have similar interests to mine.} & \enquote{Leute… die ähnliche Interessen haben wie ich.} \\
I4 & \enquote{Hmmm.} & \enquote{Hmmm.} \\
I4 & \enquote{what the tool suggests to me.} & \enquote{was das Tool mir vorschlägt.} \\
I4 & \enquote{I want to have the freedom as well.} & \enquote{möchte ich auch die Freiheit haben.} \\
I4 & \enquote{Peer-based recommendations acceptable as optional / not as a main feature.} & \enquote{Peer-based recommendations acceptable as optional} / \enquote{nicht als Main Feature.} \\
I4 & \enquote{for a reference point.} & \enquote{für einen Richtwert.} \\
I4 & \enquote{that you simply don't display it.} & \enquote{dass man es halt net anzeigt.} \\
I5 & \enquote{I took most of it from the recommendation.} & \enquote{das meiste habe ich aus der Recommendation genommen.} \\
I5 & \enquote{about the effort of the courses.} & \enquote{auf den Aufwand von den LVAs.} \\
I5 & \enquote{experience reports.} & \enquote{Erfahrungsberichte.} \\
I6 & \enquote{how well it goes down with the students.} & \enquote{wie gut bei den Studenten ankommt.} \\
I6 & \enquote{whether you then want to do it yourself.} & \enquote{ob man die selber dann auch machen will.} \\
\addlinespace
\multicolumn{3}{@{}p{12.9cm}@{}}{\textbf{T24} --- Needs recommendations based on behavioral signals (popularity, scarcity, prior courses) \quad ($n=2$; Carried forward)} \\[2pt]
I1 & \enquote{…popularity… often attended… few places… register early… only 20 places…} & \enquote{...Popularity… oft besucht… wenig Platz… melde dich früher an… nur 20 Plätze…} \\
I1 & \enquote{…on the basis of what has already been attended… `Hey, you have done these courses now… have a look at…'} & \enquote{...auf Basis von bereits besuchten… `Hey, du hast jetzt die Kurse gemacht… Schau dir… an.'} \\
I5 & \enquote{Exam attempts… a good indicator.} & \enquote{Prüfungsantritte… guter Indikator.} \\
\addlinespace
\multicolumn{3}{@{}p{12.9cm}@{}}{\textbf{T25} --- Needs social alignment / cohort visibility to reduce isolation and enable repeated peer contact \quad ($n=1$; Excluded, OOS4)} \\[2pt]
I1 & \enquote{…cool… to know… in which semester… which subject… so that I… have many fellow students with me.} & \enquote{...cool… zu wissen… in welchem Semester… welches Fach… damit ich… viele Kolleginnen habe.} \\
I1 & \enquote{I was almost alone through the whole Master's… scrambled from course to course…} & \enquote{Ich war im ganzen Master fast alleine… von Kurs zu Kurs gerangelt…} \\
I1 & \enquote{…had I known that everyone does it that way… I would have… seen the same faces again and again.} & \enquote{...hätte ich gewusst, dass das alle so machen… hätte ich… Gesichter immer wieder gesehen.} \\
\addlinespace
\multicolumn{3}{@{}p{12.9cm}@{}}{\textbf{T26} --- Needs workload/effort estimation support beyond ECTS to avoid misjudgement and regret \quad ($n=5$; Excluded, OOS4)} \\[2pt]
I1 & \enquote{…20 ECTS are not 20 ECTS… 60 ECTS… something else entirely…} & \enquote{...20 ECDS sind nicht 20 ECDS… 60 ECDS… ganz was anderes…} \\
I1 & \enquote{…if… many students used it… over the years… machine learning… reference values… true ECTS value…} & \enquote{...wenn… viele Studierende verwenden würden… über Jahre hinweg… Machine Learning… Richtwerte… wahrer ECDS-Wert…} \\
I3 & \enquote{I over- or underestimate some subjects.} & \enquote{manche Fächer überschätze oder unterschätze.} \\
I3 & \enquote{I regret so much that I didn't take that subject… Interface and Design.} & \enquote{bereue es so sehr, dass ich dieses Fach nicht genommen habe… Interface und Design.} \\
I3 & \enquote{Underestimated computer graphics… much more effort.} & \enquote{Computergrafik unterschätzt… viel mehr Aufwand.} \\
I3 & \enquote{for me that's more the problem.} & \enquote{für mich ist eher das das Problem.} \\
I3 & \enquote{didn't take Interface… I do regret that.} & \enquote{Interface nicht genommen habe… Das bereue ich schon.} \\
I4 & \enquote{how many exams… how many exercises.} & \enquote{wie viele Prüfungen… wie viele Übungen.} \\
I4 & \enquote{not the classic conversion.} & \enquote{nicht das klassische Umrechnen.} \\
I4 & \enquote{very individual.} & \enquote{sehr individuell.} \\
I5 & \enquote{how difficult the subject is.} & \enquote{wie schwierig das Fach ist.} \\
I5 & \enquote{what other students say about it.} & \enquote{was andere Studenten dazu sagen.} \\
I5 & \enquote{with six ECTS… more effort.} & \enquote{bei sechs ECTS… mehr Aufwand.} \\
I5 & \enquote{a few… that are very easy.} & \enquote{ein paar… die halt sehr leicht sind.} \\
I5 & \enquote{I'd rather have… more breathing room.} & \enquote{lieber das… mehr Breathing Room.} \\
I5 & \enquote{finding courses… very hard.} & \enquote{Kurse zu finden… sehr schwer.} \\
I5 & \enquote{make a list by difficulty.} & \enquote{Liste machen von der Schwierigkeit her.} \\
I5 & \enquote{an indicator of how heavily it's weighted.} & \enquote{Indikator, wie schwer es gewichtet ist.} \\
I6 & \enquote{like the hardest task… to cobble it together yourself.} & \enquote{wie schwierigste Aufgabe… sich das selber zusammenzuschustern.} \\
I6 & \enquote{on the basis of the content.} & \enquote{auf Basis des Contents.} \\
I6 & \enquote{IT security… carries weight.} & \enquote{IT-Security… ins Gewicht nimmt.} \\
I6 & \enquote{suggest knowledge graphs.} & \enquote{Knowledge Graphs vorschlagen.} \\
I6 & \enquote{not in the same module.} & \enquote{nicht im gleichen Modul.} \\
\addlinespace
\multicolumn{3}{@{}p{12.9cm}@{}}{\textbf{T27} --- Needs clear UI meaning of spatial layout (explain encoding; padding; vertical position) \quad ($n=1$; Carried forward)} \\[2pt]
I1 & \enquote{…the layout confuses me… why such large padding… I understand the ordering… but… why… down there…} & \enquote{...Aufteilung verwirrt… warum so ein großes Padding… die Ordnung verstehe ich… aber… warum… unten…} \\
I1 & \enquote{…explanation… is there any meaning to the vertical position? / And the answer is no.} & \enquote{...Erklärung… gibt es auch irgendeine Bedeutung zur vertikalen Position?} / \enquote{Und die Antwort ist nein.} \\
\addlinespace
\multicolumn{3}{@{}p{12.9cm}@{}}{\textbf{T28} --- Needs freedom to arrange items visually to express uncertainty and personal prioritisation \quad ($n=2$; Carried forward)} \\[2pt]
I1 & \enquote{…maybe I'd rather do it like that, because then I can place them how I want… draw in room to grow…} & \enquote{...vielleicht mache ich es… lieber, weil dann kann ich sie so hinlegen, wie ich es will… Luft nach oben einzeichnen…} \\
I1 & \enquote{…up here… definitely… down here… I'm unsure.} & \enquote{...oben… unbedingt… hier unten… unsicher bin.} \\
I1 & \enquote{…core modules… actually have to happen… right at the top… transferable skill… I wouldn't mind… rather… further down…} & \enquote{...Core-Module… muss eigentlich passieren… ganz oben… Transferrable Skill… wäre mir wurscht… eher… unten…} \\
I1 & \enquote{…helps me with my mental grouping… That does help me.} & \enquote{...hilft mir für meine mentale Gruppierung… Das hilft mir schon.} \\
I1 & \enquote{…picture board… distance… I push it right to the back, because I don't care about it.} & \enquote{...Ficture-Bord… Distanz… ziehe ich es ganz nach hinten, weil es mir wurscht ist.} \\
I4 & \enquote{by priority.} & \enquote{nach Priorität.} \\
\addlinespace
\multicolumn{3}{@{}p{12.9cm}@{}}{\textbf{T29} --- Needs grouping mechanics controllable + flexible (named groups, spanning semesters, move behaviour) \quad ($n=1$; Carried forward)} \\[2pt]
I1 & \enquote{And how does this grouping come about?… can I… move it… do you do that automatically?} & \enquote{Und wie entsteht diese Gruppierung?… kann ich… verschieben… macht ihr das von alleine?} \\
I1 & \enquote{…if I've the exercise in the first… in the third… does it then span across? / …maybe you would want to shape that even more freely.} & \enquote{...wenn ich die Übung im ersten… im dritten… geht das dann so drüber?} / \enquote{...vielleicht will man das gerne noch freier gestalten.} \\
I1 & \enquote{…the possibility to freely create groups… and name them how I want.} & \enquote{...Möglichkeit… frei… Gruppierungen erstelle… benennen kann, wie ich sie will.} \\
\addlinespace
\multicolumn{3}{@{}p{12.9cm}@{}}{\textbf{T30} --- Needs to express soft/implicit dependencies and decide between grouping vs edges \quad ($n=2$; Carried forward)} \\[2pt]
I1 & \enquote{…dependencies… good if I… only do it once I've the first three… maybe… draw them freely.} & \enquote{...Abhängigkeiten… gut, wenn ich… erst dann mach, wenn ich die ersten drei hab… vielleicht… frei zeichnen.} \\
I1 & \enquote{…it isn't stated explicitly anywhere, we don't know that, of course.} & \enquote{...da steht nirgends explizit, das wissen wir natürlich nicht.} \\
I1 & \enquote{I wouldn't know what for.} & \enquote{Wüsste ich nicht, wozu.} \\
I2 & \enquote{not compulsory… but it would be good to do that one first.} & \enquote{keine Pflicht… aber es wäre gut, dass du das vor dem machst.} \\
\addlinespace
\multicolumn{3}{@{}p{12.9cm}@{}}{\textbf{T31} --- Needs prerequisite/dependency representation (formal + recommended) and order enforcement/warnings \quad ($n=6$; Carried forward)} \\[2pt]
I1 & \enquote{I think that's good, yes.} & \enquote{Das finde ich gut, ja.} \\
I1 & \enquote{But not… building on each other, right? / No, simply similar content.} & \enquote{Aber nicht… aufbauend, ne?} / \enquote{Nein, einfach ähnliche Inhalte.} \\
I2 & \enquote{There is no subject that rules you out completely… They are only recommendations.} & \enquote{Es gibt kein Fach, was dich komplett ausschließt… Es sind nur Empfehlungen.} \\
I2 & \enquote{EP1 and EP2… they directly build on each other.} & \enquote{EP1 und EP2… die bauen halt direkt davon aufeinander.} \\
I2 & \enquote{I'm missing… the dependencies between the subjects.} & \enquote{mir fehlen… die Dependencies zwischen den Fächern.} \\
I2 & \enquote{that the system… warns… or doesn't allow it at all.} & \enquote{dass das System… warnt… oder gar nicht erlaubt.} \\
I2 & \enquote{you can't… do one and two… in the same semester.} & \enquote{man kann… eins und zwei… nicht an einem Semester machen.} \\
I3 & \enquote{I read through what you need to be able to do… these prerequisites.} & \enquote{lese ich mir halt durch, was man alles können muss… diese Vorkenntnisse.} \\
I4 & \enquote{how they are connected to each other… dependencies between courses.} & \enquote{wie die miteinander verbunden sind… Dependencies zwischen Kurse.} \\
I4 & \enquote{it makes sense to do the first stage first.} & \enquote{macht's ja Sinn, dass man zuerst die erste Stufe macht.} \\
I4 & \enquote{you have to have done it… not before it in time.} & \enquote{man muss ihn gemacht haben… nicht zeitlich davor.} \\
I4 & \enquote{Basic edges.} & \enquote{Basic-Kanten.} \\
I4 & \enquote{highlight more.} & \enquote{mehr hervorheben.} \\
I4 & \enquote{by colour.} & \enquote{farblich.} \\
I5 & \enquote{normally… one first and then.} & \enquote{im Normalfall… zuerst eins und dann.} \\
I6 & \enquote{one of the best representations.} & \enquote{eine der besten Darstellungen.} \\
I6 & \enquote{very difficult to make out.} & \enquote{sehr schwierig zu erkennen.} \\
I6 & \enquote{so that… I don't get into trouble.} & \enquote{damit… nicht in die Bredouille komme.} \\
\addlinespace
\multicolumn{3}{@{}p{12.9cm}@{}}{\textbf{T32} --- Needs graph layout consistent with reading direction or alternatively full rearrangement freedom \quad ($n=1$; Carried forward)} \\[2pt]
I1 & \enquote{…nicer if it goes left to right… because we read left to right.} & \enquote{...schöner, wenn man von links nach rechts geht… weil man liest ja von links nach rechts.} \\
I1 & \enquote{So either that, or I've the freedom to do it myself.} & \enquote{Also entweder das oder ich habe halt Freiheiten.} \\
\addlinespace
\multicolumn{3}{@{}p{12.9cm}@{}}{\textbf{T33} --- Needs tool to be tailored to the specific programme structure to be adopted and trusted \quad ($n=2$; Carried forward)} \\[2pt]
I1 & \enquote{…very tailored to Software Engineering… definitely.} & \enquote{...sehr maßgeschneidert für Software Engineering… auf jeden Fall.} \\
I1 & \enquote{…the core modules are in there… mandatory modules… I can rely on that.} & \enquote{...da sind die Core Module drin… Mandatory Module… da kann ich mich… verlassen.} \\
I4 & \enquote{cool that it's already listed up front.} & \enquote{cool, dass das schon vor Haus aufgelistet ist.} \\
I4 & \enquote{I don't have to look it up.} & \enquote{Ich muss nicht raussuchen.} \\
\addlinespace
\multicolumn{3}{@{}p{12.9cm}@{}}{\textbf{T34} --- Needs possibility to continue collaboration / participate in follow-up \quad ($n=4$; Methodological artefact, asking participants about a follow-up round)} \\[2pt]
I1 & \enquote{Yes… follow-up… Then I can bring my Excel along too.} & \enquote{Ja… Follow-up… Dann kann ich auch mein Excel mitbringen.} \\
I2 & \enquote{Yes, gladly.} & \enquote{Ja, gerne.} \\
I3 & \enquote{Yes, yes.} & \enquote{Ja, ja.} \\
I4 & \enquote{Okay, I see.} & \enquote{Okay, verstehe.} \\
I4 & \enquote{Yes.} & \enquote{Ja.} \\
\addlinespace
\multicolumn{3}{@{}p{12.9cm}@{}}{\textbf{T35} --- Needs semester capacity targeting (e.g., aim for ~30~ECTS) and capacity-based planning \quad ($n=4$; Carried forward)} \\[2pt]
I2 & \enquote{so that I'm registered for around thirty.} & \enquote{damit ich halt quasi um die… dreißig angemeldet bin.} \\
I2 & \enquote{One, one ECTS you can somehow pick up with optional things…} & \enquote{Ein, ein ECTS kann man sich irgendwie optionale Sachen holen…} \\
I2 & \enquote{fifty ECTS… written in… collisions came up… dropped down to thirty.} & \enquote{fünfzig ECTS… reingeschrieben… ergaben sich Kollisionen… auf dreißig runtergesprungen.} \\
I2 & \enquote{I've… twenty-seven… still need three.} & \enquote{ich habe… siebenundzwanzig… brauche noch drei.} \\
I3 & \enquote{how many ECTS the individual modules… have.} & \enquote{wie viele ECTS die einzelnen Module… haben.} \\
I5 & \enquote{to… see whether the pace suits me.} & \enquote{um… zu schauen, ob der Pace für mich passt.} \\
I5 & \enquote{31~ECTS… that's too much.} & \enquote{31 ECTS… das ist zu viel.} \\
I5 & \enquote{brought it down… to about 21~ECTS.} & \enquote{runtergebraucht… auf 21 ECTS ungefähr.} \\
I6 & \enquote{how many ECTS you do in the semesters.} & \enquote{wie viel ECTS man in den Semestern macht.} \\
I6 & \enquote{how many ECTS I do in a semester.} & \enquote{wie viel ECTS ich in einem Semester mache.} \\
I6 & \enquote{that I do enough ECTS.} & \enquote{dass ich genügend ECTS mache.} \\
I6 & \enquote{overview… of my ECTS.} & \enquote{Überblick… über meine ECTS.} \\
\addlinespace
\multicolumn{3}{@{}p{12.9cm}@{}}{\textbf{T36} --- Needs semester collision detection and early exam schedule availability \quad ($n=1$; Excluded, OOS3)} \\[2pt]
I2 & \enquote{that exams don't overlap.} & \enquote{dass Prüfungen sich nicht überschneiden.} \\
I2 & \enquote{the exams… not written down on the first day of term.} & \enquote{die Prüfungen… nicht aufgeschrieben am ersten Unitag.} \\
I2 & \enquote{That didn't work in the third one.} & \enquote{Das hat dann im Dritten nicht funktioniert.} \\
I2 & \enquote{you have to build a table and enter all the dates yourself.} & \enquote{man muss sich Tabelle herstellen und die ganzen Termine sich eintragen.} \\
I2 & \enquote{Excel sheet… where I enter all the dates.} & \enquote{Excel-Tabelle… wo ich die ganzen Termine eintrage.} \\
I2 & \enquote{It took me three days this semester… you write down all the dates.} & \enquote{Drei Tage hat es mir dieses Semester gedauert… schreibt man die ganzen Termine runter.} \\
I2 & \enquote{after three weeks it's: `Okay, we are going to run the exam completely differently now.'} & \enquote{nach drei Wochen heißt es: `Okay, wir schreiben jetzt mal ganz anders die Prüfung.'} \\
I2 & \enquote{Key feature… semester view… collisions.} & \enquote{Key Feature… Semesteransicht… Kollisionen.} \\
I2 & \enquote{that I no longer need this Excel sheet.} & \enquote{dass mir diese Excel-Tabelle nicht mehr nötig ist.} \\
I2 & \enquote{vertically… zoom into the semesters… scheduling conflicts.} & \enquote{vertikal… in die Semester reinzoomt… Zeitliche Konflikte.} \\
I2 & \enquote{that's then not the system's problem.} & \enquote{das ist dann nicht das Problem des Systems.} \\
\addlinespace
\multicolumn{3}{@{}p{12.9cm}@{}}{\textbf{T37} --- Needs to reduce manual effort/time in planning (data entry, repetitive checking) \quad ($n=3$; Carried forward)} \\[2pt]
I2 & \enquote{checks it again in the… course text.} & \enquote{prüft dann noch mal im… LVUA-Text.} \\
I2 & \enquote{it's time-consuming… definitely.} & \enquote{es ist zeitaufwendig… auf jeden Fall.} \\
I2 & \enquote{your way… is simply to put more time in.} & \enquote{deine Art… ist einfach mehr Zeit reinstecken.} \\
I4 & \enquote{that the plan already takes quite a bit off your shoulders.} & \enquote{dass der Plan da schon einiges abnimmt.} \\
I6 & \enquote{clicking through TISS… pulling it out manually.} & \enquote{durch TIS durchklicken… manuell rausholen.} \\
I6 & \enquote{not really a nice planning workflow.} & \enquote{nicht wirklich ein schöner Planning-Workflow.} \\
\addlinespace
\multicolumn{3}{@{}p{12.9cm}@{}}{\textbf{T38} --- Needs planning that reflects actual course availability over time (offerings/timing constraints) \quad ($n=5$; Excluded, OOS3)} \\[2pt]
I2 & \enquote{Digital Health specialisation… but the subjects weren't offered… not at all in the first three semesters.} & \enquote{Digital-Health-Vertiefung… aber die Fächer wurden nicht angeboten… die ersten drei Semester gar nicht.} \\
I2 & \enquote{Because nothing was offered… Now they are slowly coming in.} & \enquote{Weil nichts angeboten wurde… Jetzt kommen die langsam rein.} \\
I3 & \enquote{I look at TISS… what is on offer in the winter semester.} & \enquote{ich guck auf TIS… was es eben im Wintersemester gibt.} \\
I4 & \enquote{not all of them offered in every semester.} & \enquote{nicht alle immer in jedem Semester angeboten.} \\
I4 & \enquote{wasn't offered… shortly before the start of the semester.} & \enquote{nicht angeboten worden… kurz vor Semesterbeginn.} \\
I4 & \enquote{tedious to find replacement courses.} & \enquote{zäh, Ersatzlehrveranstaltungen zu finden.} \\
I5 & \enquote{took whatever was available at the time.} & \enquote{genommen, was gerade zur Verfügung war.} \\
I5 & \enquote{some are only offered in the winter semester.} & \enquote{manche gibt es ja nur im Wintersemester.} \\
I5 & \enquote{that was often the biggest problem.} & \enquote{das war das größte Problem oft.} \\
I5 & \enquote{with my friend… let us take this course now.} & \enquote{mit meinem Freund… lass jetzt diese LVA nehmen.} \\
I6 & \enquote{courses cancelled… postponed.} & \enquote{Lehrveranstaltungen abgesagt… verschoben.} \\
\addlinespace
\multicolumn{3}{@{}p{12.9cm}@{}}{\textbf{T39} --- Needs drill-down access to official course information from the plan/graph \quad ($n=3$; Excluded, OOS3)} \\[2pt]
I2 & \enquote{if you click on the little block… TISS page… popup.} & \enquote{wenn man auf das Blöckchen drauf klickt… TIS-Seite… Popup.} \\
I4 & \enquote{you click on it and then it pops up.} & \enquote{raufklickst und dann ploppt es auf.} \\
I6 & \enquote{Link… to the course page.} & \enquote{Link… zur Vobi-Seite.} \\
\addlinespace
\multicolumn{3}{@{}p{12.9cm}@{}}{\textbf{T40} --- Needs intuitive labels/legend for course types and symbols \quad ($n=1$; Carried forward)} \\[2pt]
I2 & \enquote{the letters not intuitively understandable.} & \enquote{die Buchstaben nicht intuitiv nachvollziehbar.} \\
\addlinespace
\multicolumn{3}{@{}p{12.9cm}@{}}{\textbf{T41} --- Needs graph purpose/interpretability clarity (what it's for, how to read it) \quad ($n=2$; Carried forward)} \\[2pt]
I2 & \enquote{at first glance I wouldn't know.} & \enquote{auf den ersten Blick wüsste ich es nicht.} \\
I2 & \enquote{nice visualisation… can't see… the purpose.} & \enquote{schöne Visualisierung… den Zweck… nicht erkennen.} \\
I4 & \enquote{what do these stand for?} & \enquote{wofür stehen die?} \\
I4 & \enquote{Initial confusion due to insufficient visual cues.} & \enquote{Initial confusion due to insufficient visual cues} \\
\addlinespace
\multicolumn{3}{@{}p{12.9cm}@{}}{\textbf{T42} --- Needs filters to reduce overload and support targeted browsing (negated + type/ECTS filters; hide/show) \quad ($n=4$; Carried forward)} \\[2pt]
I2 & \enquote{Filter… that's way too much at the beginning.} & \enquote{Filter… das ist viel zu viel am Anfang.} \\
I2 & \enquote{I would like… to see `not added courses'.} & \enquote{ich möchte halt… `not added courses' sehen.} \\
I2 & \enquote{hide and show transferable skills.} & \enquote{transable Skills quasi ausblenden, einblenden.} \\
I2 & \enquote{Slider… between three and five.} & \enquote{Schieberegler… zwischen drei und fünf.} \\
I3 & \enquote{so that it isn't so much.} & \enquote{damit es nicht so viel ist.} \\
I3 & \enquote{mandatory… narrow elective… broad elective.} & \enquote{Pflichtfach… enge Wahl… breite Wahl.} \\
I3 & \enquote{what interests me in principle.} & \enquote{was mich interessiert grundsätzlich.} \\
I3 & \enquote{I would generally always want to see the compulsory subjects.} & \enquote{Pflichtfächer würde ich generell immer sehen wollen.} \\
I3 & \enquote{broad elective… collapse and expand.} & \enquote{breite Wahl… ein- und aufklappen.} \\
I3 & \enquote{if even more modules come in, then it's an extreme amount.} & \enquote{wenn dann noch mehr Module kommen, dann ist es extrem viel.} \\
I4 & \enquote{set… which semester.} & \enquote{einstellen… welches Semester.} \\
I4 & \enquote{I would like a filter for that.} & \enquote{da hätte ich gerne einen Filter.} \\
I4 & \enquote{Semester.} & \enquote{Semester.} \\
I4 & \enquote{only the recommended ones.} & \enquote{nur die empfohlenen.} \\
I4 & \enquote{which ones I've already completed.} & \enquote{welche ich schon absolviert habe.} \\
I4 & \enquote{only three ECTS.} & \enquote{nur drei ECTS.} \\
I6 & \enquote{which ones I've completed.} & \enquote{welche ich abgeschlossen habe.} \\
I6 & \enquote{collapse.} & \enquote{einklappen.} \\
I6 & \enquote{only the items still open.} & \enquote{nur mehr die Items offen.} \\
\addlinespace
\multicolumn{3}{@{}p{12.9cm}@{}}{\textbf{T43} --- Needs recommendation UI that annotates rather than hides, and keeps layout stable \quad ($n=2$; Carried forward)} \\[2pt]
I2 & \enquote{marked… but not destroying the graph.} & \enquote{markiert… aber nicht den Graph zerstört.} \\
I2 & \enquote{Filter… removing things… suggestions… only marking.} & \enquote{Filter… Sachen entfernen… Vorschläge… nur markieren.} \\
I4 & \enquote{glowing borders.} & \enquote{erleuchtende Umrandungen.} \\
I4 & \enquote{the ones I've already completed… on there as well.} & \enquote{die ich schon absolviert habe… auch drauf.} \\
\addlinespace
\multicolumn{3}{@{}p{12.9cm}@{}}{\textbf{T44} --- Needs support to move from overview graph to actionable semester plan \quad ($n=3$; Carried forward)} \\[2pt]
I2 & \enquote{what would you need… to get from your graph… to planning?} & \enquote{was bräuchtest du… dass du von deinem Graph kommst… zum Planen?} \\
I3 & \enquote{Add to plan… I choose my semester.} & \enquote{Add to plan… wähle mir mein Semester aus.} \\
I4 & \enquote{Popup… add.} & \enquote{Popup… hinzufügen.} \\
\addlinespace
\multicolumn{3}{@{}p{12.9cm}@{}}{\textbf{T45} --- Needs recommendations based on interests + content similarity + dependency relations \quad ($n=4$; Carried forward)} \\[2pt]
I2 & \enquote{Software Engineering Project builds indirectly on… Web Engineering.} & \enquote{Software Engineering Projekt basiert indirekt auf… Web Engineering.} \\
I2 & \enquote{other subjects that cover similar things.} & \enquote{andere Fächer, die ähnliche Sachen abdecken.} \\
I2 & \enquote{similar to this suggestion function.} & \enquote{ähnlich wie diese Vorschlagsfunktion.} \\
I2 & \enquote{would definitely be nice.} & \enquote{wäre auf jeden Fall nett.} \\
I3 & \enquote{thematic recommendations… that's what I'm working towards as well.} & \enquote{thematische Empfehlungen… darauf arbeite ich auch hin.} \\
I3 & \enquote{courses that match my interests.} & \enquote{Kurse, die zu meinen Interessen passen.} \\
I3 & \enquote{similar courses… that would be helpful for you.} & \enquote{ähnliche Kurse… das wäre für dich hilfreich.} \\
I4 & \enquote{By interest, basically.} & \enquote{Nach Interesse halt.} \\
I4 & \enquote{the ones that are connected.} & \enquote{die zusammenhängen.} \\
I4 & \enquote{similar… in terms of content.} & \enquote{ähnliche… inhaltmäßig.} \\
I4 & \enquote{I think that fits well.} & \enquote{passt finde ich gut.} \\
I6 & \enquote{reflects… my preferences.} & \enquote{meine Präferenzen… abbildet.} \\
I6 & \enquote{on the basis of the content.} & \enquote{auf Basis des Contents.} \\
I6 & \enquote{IT security… carries weight.} & \enquote{IT-Security… ins Gewicht nimmt.} \\
I6 & \enquote{suggest knowledge graphs.} & \enquote{Knowledge Graphs vorschlagen.} \\
I6 & \enquote{not in the same module.} & \enquote{nicht im gleichen Modul.} \\
\addlinespace
\multicolumn{3}{@{}p{12.9cm}@{}}{\textbf{T46} --- Needs planning that is robust to schedule instability (boundary of what system can/can't solve) \quad ($n=1$; Carried forward)} \\[2pt]
I2 & \enquote{after three weeks it's: `Okay, we are going to run the exam completely differently now.'} & \enquote{nach drei Wochen heißt es: `Okay, wir schreiben jetzt mal ganz anders die Prüfung.'} \\
I2 & \enquote{that's then not the system's problem.} & \enquote{das ist dann nicht das Problem des Systems.} \\
\addlinespace
\multicolumn{3}{@{}p{12.9cm}@{}}{\textbf{T47} --- Needs a table-based semester planning view that supports prioritisation, notes, and preferred layout \quad ($n=2$; Carried forward)} \\[2pt]
I3 & \enquote{that you have a plan… what you study… that you set priorities, which subject has priority and which doesn't.} & \enquote{dass man einen Plan hat… was man lernt… sich Prioritäten setzt, welches Fach hat Priorität und welches nicht.} \\
I3 & \enquote{I simply find it easier with tables.} & \enquote{Ich tue mir mit Tabellen einfach leichter.} \\
I3 & \enquote{I think that should probably fit.} & \enquote{ich glaube, das sollte so vielleicht passen.} \\
I3 & \enquote{that I've space… to write something next to it.} & \enquote{Platz habe… etwas daneben hinzuschreiben.} \\
I3 & \enquote{which subjects have priority.} & \enquote{welche Fächer haben Priorität.} \\
I3 & \enquote{only one exam at the end… lower priority.} & \enquote{nur eine Prüfung am Ende… weniger Priorität.} \\
I3 & \enquote{Algorithms and Data Structures… I have to work on it continuously through the semester.} & \enquote{AlgoDat… muss ich im Semester durchgehend machen.} \\
I3 & \enquote{that wouldn't really help me now.} & \enquote{mir würde das jetzt nicht wirklich helfen.} \\
I3 & \enquote{Yes… the table, definitely.} & \enquote{Ja… die Tabelle auf jeden Fall.} \\
I5 & \enquote{looks very clear.} & \enquote{seh sehr übersichtlich.} \\
\addlinespace
\multicolumn{3}{@{}p{12.9cm}@{}}{\textbf{T48} --- Needs competency- and learning-outcome-based planning support (self-assessment, reflection, prerequisite readiness) \quad ($n=2$; Carried forward)} \\[2pt]
I3 & \enquote{What can you do? What can you not do?… if you can't program well, then… maybe rather not take the programming-heavy subjects.} & \enquote{Was kann man? Was kann man nicht?… wenn man nicht gut programmieren kann, dann… programmierschwere Fächer eher vielleicht nicht nehmen.} \\
I3 & \enquote{theoretical computer science… I already have prior knowledge there.} & \enquote{theoretische Informatik… habe ich schon Vorwissen.} \\
I3 & \enquote{what did I learn from these subjects? And then you plan… the next semester.} & \enquote{was habe ich aus diesen Fächern gelernt? Und dann plant man… das nächste Semester.} \\
I5 & \enquote{I was… able to program.} & \enquote{ich war… kann programmieren.} \\
\addlinespace
\multicolumn{3}{@{}p{12.9cm}@{}}{\textbf{T49} --- Needs support for balancing assessment types and distributing workload/exam load across the year \quad ($n=2$; Carried forward)} \\[2pt]
I3 & \enquote{two subjects… project… no exam… and then… subjects… where you do have an exam.} & \enquote{zwei Fächer… Projekt… keine Prüfung… und dann… Fächer… wo man eben eine Prüfung hat.} \\
I3 & \enquote{spread the workload a little across the semester.} & \enquote{Workload ein bisschen aufteilt über Semester.} \\
I3 & \enquote{four exams in January… for me that would be too much.} & \enquote{vier Prüfungen im Jänner… für mich wäre das zu viel.} \\
I5 & \enquote{not too many extremely difficult courses per semester.} & \enquote{nicht zu viele extrem schwierige LVAs pro Semester.} \\
I5 & \enquote{balanced.} & \enquote{ausbalanciert.} \\
I5 & \enquote{seen from experience.} & \enquote{aus Erfahrung gesehen.} \\
I5 & \enquote{have the difficulty balanced.} & \enquote{Schwierigkeit ausbalanciert haben.} \\
I5 & \enquote{not a good strategy.} & \enquote{nicht 'ne gute Strategie.} \\
I5 & \enquote{Semester… one of five.} & \enquote{Semester… eins von fünf.} \\
\addlinespace
\multicolumn{3}{@{}p{12.9cm}@{}}{\textbf{T50} --- Needs scalable visualisation via abstraction (avoid graph \enquote{explosion}; show summaries, keep details elsewhere; graph optional) \quad ($n=1$; Carried forward)} \\[2pt]
I3 & \enquote{I would maybe… only write down Computer Systems.} & \enquote{ich würd vielleicht… nur Computersysteme… hinschreiben.} \\
I3 & \enquote{you do that separately in a table.} & \enquote{das macht man separat in einer Tabelle.} \\
I3 & \enquote{with graphs, it kind of explodes.} & \enquote{bei Grafen explodiert das irgendwie.} \\
I3 & \enquote{good to have in addition.} & \enquote{zusätzlich gut zu haben.} \\
I3 & \enquote{nice to have.} & \enquote{nice to have.} \\
\addlinespace
\multicolumn{3}{@{}p{12.9cm}@{}}{\textbf{T51} --- Needs lightweight content clarification beyond official descriptions (external lookup) \quad ($n=1$; Excluded, OOS3)} \\[2pt]
I3 & \enquote{Maybe I google something to see what it might be about.} & \enquote{Vielleicht google ich etwas, um zu schauen, um was es vielleicht gehen wird.} \\
\addlinespace
\multicolumn{3}{@{}p{12.9cm}@{}}{\textbf{T52} --- Needs career/internship-aligned planning option (course choices toward internship readiness) \quad ($n=1$; Carried forward)} \\[2pt]
I3 & \enquote{which subjects do I need… in order to do an internship?} & \enquote{welche Fächer brauche ich… um ein Praktikum zu machen?} \\
I3 & \enquote{not a must.} & \enquote{nicht ein Must.} \\
\addlinespace
\multicolumn{3}{@{}p{12.9cm}@{}}{\textbf{T53} --- Needs frictionless, frequent access to the plan (mobile/ambient visibility) to enable real usage \quad ($n=1$; Excluded, desktop-scale web application (Section~\ref{sec:overall-scope}))} \\[2pt]
I4 & \enquote{didn't have it open regularly.} & \enquote{nicht regelmäßig offen gehabt.} \\
I4 & \enquote{that only works on the laptop.} & \enquote{das geht nur am Laptop.} \\
I4 & \enquote{an app… background image… plug-in.} & \enquote{eine App… Hintergrundbild… Plug-in.} \\
I4 & \enquote{like a calendar, basically.} & \enquote{wie ein Kalender halt.} \\
I4 & \enquote{Logging in… a little cumbersome.} & \enquote{Einloggen… ein bisschen umständlich.} \\
I4 & \enquote{Widget… on the phone.} & \enquote{Widget… am Handy.} \\
I4 & \enquote{What good is it to me if I've the plan but never look at it.} & \enquote{Was bringt mir das, wenn ich den Plan hab, aber ich schau den nie an.} \\
\addlinespace
\multicolumn{3}{@{}p{12.9cm}@{}}{\textbf{T54} --- Needs visual-first, low-text UI (icons, colour, animation, aesthetic polish) to support motivation and comprehension \quad ($n=2$; Carried forward)} \\[2pt]
I4 & \enquote{with icons… not much text.} & \enquote{mit Icons… nicht viel Text.} \\
I4 & \enquote{too much text… off-putting.} & \enquote{zu viel Text… abschreckend.} \\
I4 & \enquote{sheet of paper… for exams.} & \enquote{Blatt Papier… für Prüfungen.} \\
I4 & \enquote{coloured in… [the source records this excerpt under T19 as well]} & \enquote{eingefärbt…} \\
I4 & \enquote{I'm… more of a visual person.} & \enquote{ich bin… visuell eher.} \\
I4 & \enquote{a bit of animation.} & \enquote{ein bisschen Animation.} \\
I4 & \enquote{visually very appealing… feel like doing it.} & \enquote{optisch sehr ansprechend… Bock haben.} \\
I4 & \enquote{bring in a bit more visually as well.} & \enquote{visuell halt auch ein bisschen mehr reinbringen.} \\
I5 & \enquote{I understand the colour coding.} & \enquote{das Color Coding verstehe ich.} \\
\addlinespace
\multicolumn{3}{@{}p{12.9cm}@{}}{\textbf{T55} --- Needs control of recommendation presentation (separation of spaces, avoid clutter; keep context) \quad ($n=2$; Carried forward)} \\[2pt]
I4 & \enquote{otherwise it's too much on one page.} & \enquote{sonst zu viel auf einer Page.} \\
I4 & \enquote{Recommendation separation to avoid clutter.} & \enquote{Recommendation separation to avoid clutter} \\
I4 & \enquote{the ones I've already completed… on there as well. [the source records this excerpt under T43 as well]} & \enquote{die ich schon absolviert habe… auch drauf.} \\
I4 & \enquote{Filter-based recommendation display preferred / it filters everything away.} & \enquote{Filter-based recommendation display preferred} / \enquote{filtert alles weg.} \\
I4 & \enquote{Filter-only recommendations perceived as dry / too dry.} & \enquote{Filter-only recommendations perceived as dry} / \enquote{zu trocken.} \\
I6 & \enquote{not here… rather in the smaller area.} & \enquote{nicht hier… eher in dem kleineren Bereich.} \\
I6 & \enquote{Subsection… recommendations.} & \enquote{Subsection… Empfehlungen.} \\
\addlinespace
\multicolumn{3}{@{}p{12.9cm}@{}}{\textbf{T56} --- Needs support for quick replanning when external disruptions happen (last-minute offering changes) \quad ($n=3$; Carried forward)} \\[2pt]
I4 & \enquote{wasn't offered… shortly before the start of the semester.} & \enquote{nicht angeboten worden… kurz vor Semesterbeginn.} \\
I4 & \enquote{tedious to find replacement courses.} & \enquote{zäh, Ersatzlehrveranstaltungen zu finden.} \\
I5 & \enquote{simply to take a different subject.} & \enquote{einfach ein anderes Fach zu nehmen.} \\
I6 & \enquote{courses cancelled… postponed.} & \enquote{Lehrveranstaltungen abgesagt… verschoben.} \\
\addlinespace
\multicolumn{3}{@{}p{12.9cm}@{}}{\textbf{T57} --- Needs module-based overview representation as a planning mental model (modules as primary units) \quad ($n=2$; Carried forward)} \\[2pt]
I4 & \enquote{per module a bar.} & \enquote{je Modul halt eine Säule.} \\
I4 & \enquote{one sheet… all the courses… another sheet for the overview.} & \enquote{ein Blatt… alle Lehrveranstaltungen… ein anderes Blatt zur Übersicht.} \\
I5 & \enquote{that you see directly… these actually belong together.} & \enquote{dass man direkt sieht… die gehören eigentlich zusammen.} \\
I5 & \enquote{it pays off to do both at the same time.} & \enquote{zahlt es sich aus, beide gleichzeitig zu machen.} \\
I5 & \enquote{it doesn't help much… next semester.} & \enquote{bringt ja nicht viel… nächstes Semester.} \\
\addlinespace
\multicolumn{3}{@{}p{12.9cm}@{}}{\textbf{T58} --- Needs explicit \enquote{planning burden reduction} (system should take over cognitive load) \quad ($n=1$; Carried forward)} \\[2pt]
I4 & \enquote{that the plan already takes quite a bit off your shoulders.} & \enquote{dass der Plan da schon einiges abnimmt.} \\
\addlinespace
\multicolumn{3}{@{}p{12.9cm}@{}}{\textbf{T59} --- Needs onboarding/support for the planning transition when entering university (planning \enquote{shock}) \quad ($n=2$; Carried forward)} \\[2pt]
I5 & \enquote{with the whole change from secondary school to university… completely new to have to plan everything about it.} & \enquote{bei der ganzen Umstellung von HTL auf Uni… ganz neu, sich alles darüber planen zu müssen.} \\
I6 & \enquote{for first-years… one of the most difficult things.} & \enquote{für Erstis… eines der schwierigsten Dinge.} \\
I6 & \enquote{it overwhelms you right at the start.} & \enquote{überfordert einen schon zu Beginn.} \\
\addlinespace
\multicolumn{3}{@{}p{12.9cm}@{}}{\textbf{T60} --- Needs risk-aware planning support (buffers, failure-risk awareness, conservative planning) \quad ($n=1$; Excluded, OOS1)} \\[2pt]
I5 & \enquote{I'm a risk-averse person.} & \enquote{ich bin eine risk-aversive Person.} \\
I5 & \enquote{if I… don't pass… ECTS lost.} & \enquote{wenn ich… nicht bestehe… ECTS verloren.} \\
\addlinespace
\multicolumn{3}{@{}p{12.9cm}@{}}{\textbf{T61} --- Needs configurable ordering + auto-alignment behaviours in planning views (sort logic, snap-to-order) \quad ($n=1$; Carried forward)} \\[2pt]
I5 & \enquote{structured from top to bottom.} & \enquote{strukturiert von oben nach unten.} \\
I5 & \enquote{that it jumps back to the top.} & \enquote{dass es sich nach oben springt.} \\
I5 & \enquote{reshuffle… up, down.} & \enquote{umschichten… nach oben, nach unten.} \\
I5 & \enquote{some… alphabetically.} & \enquote{manche… alphabetisch.} \\
I5 & \enquote{by difficulty.} & \enquote{nach Schwierigkeiten.} \\
\addlinespace
\multicolumn{3}{@{}p{12.9cm}@{}}{\textbf{T62} --- Needs automation/data integration to reduce errors from manual spreadsheet workflows \quad ($n=1$; Carried forward)} \\[2pt]
I6 & \enquote{information… pulled in… automatically.} & \enquote{Informationen… automatisch… gezogen.} \\
I6 & \enquote{would simplify the whole workflow.} & \enquote{würde den ganzen Workflow vereinfachen.} \\
I6 & \enquote{very error-prone… manual.} & \enquote{sehr fehleranfällig… manuell.} \\
I6 & \enquote{you might slip in the row.} & \enquote{man verrutscht vielleicht in der Zeile.} \\
I6 & \enquote{no Excel expert… pretty basic.} & \enquote{kein Excel-Experte… pretty basic.} \\
\addlinespace
\multicolumn{3}{@{}p{12.9cm}@{}}{\textbf{T63} --- Needs recommendation presentation that highlights institutional \enquote{recommended curriculum} (explicit labelling/trust cue) \quad ($n=1$; Carried forward)} \\[2pt]
I6 & \enquote{recommended by the… institute.} & \enquote{empfohlen wird vom… Institut.} \\
I6 & \enquote{really marked as such.} & \enquote{wirklich gekennzeichnet.} \\
I6 & \enquote{recommended curriculum.} & \enquote{empfohlenen Curriculum.} \\
\addlinespace
\multicolumn{3}{@{}p{12.9cm}@{}}{\textbf{T64} --- Needs an overload-safe course catalogue experience (avoid \enquote{dumping everything}; relevance-filtered suggestions) \quad ($n=1$; Carried forward)} \\[2pt]
I6 & \enquote{an extreme number of courses thrown at you.} & \enquote{extrem viele Lehrveranstaltungen hingeworfen.} \\
I6 & \enquote{courses… suggested… relevant.} & \enquote{Lehrveranstaltungen… vorgeschlagen… relevant.} \\
I6 & \enquote{you lose track very quickly.} & \enquote{sehr schnell den Überblick verliert.} \\
\addlinespace
\end{longtable}
\normalsize

\section{Feature Specifications}
\label{app:feature-specs}
This appendix records in full the 13 features derived in Section~\ref{chap:from-themes-to-requirements}, which states the derivation procedure and summarises the result in Table~\ref{tab:feature-summary}. They are reproduced here so that every requirement cited from the Design and Implementation chapters can be read in its own evidential context. Features appear in identifier order, which Table~\ref{tab:feature-summary} reorders by importance score, and each is presented in the same structure: title and identifier, importance score, associated themes, description and user evidence, derived requirements, and a delineation where one applies. The themes named under each feature are defined, with the excerpts that informed them, in Section~\ref{app:codebook}; the 13 features carry 55 numbered requirements between them.

\newcounter{featurenum}

\subsection{FEATURE-001 --- Unified Data \& Tool Integration}
\refstepcounter{featurenum}\label{feat:001}
\textbf{Importance score 20.} Associated themes: T9, T10, T11, T12, T15, T37, T62.

\paragraph{Description and user evidence.}
Participants plan through a self-assembled tool ecology rather than one system, combining spreadsheets, notes and the institutional platform TISS (T10), which is itself described as unclear and computationally passive: \enquote{the overview isn't clear… it doesn't calculate anything} (T9). Transferring course data out of it by hand is burdensome and error-prone, \enquote{very error-prone… manual… you might slip in the row} (T62), which motivates automatic import. The ecology is sustained by informal peer reuse, a friend's spreadsheet passed on because \enquote{no finished Excel from the university} exists (T11), and participants want a better representation than TISS offers, with personal tracking fields and visualisation (T12, T15, T37).

\paragraph{Requirements.}
\begin{enumerate}
    \item The system can initialise a plan without manual course data entry, by importing at minimum: programme/module structure, course list, ECTS, and planning-relevant course metadata.
    \item The system maintains stable internal identifiers for courses/modules such that user selections persist across updates and across table/graph views.
    \item Users can create a semester plan without copying course data manually.
    \item The system supports user-owned planning fields (at minimum: notes, grades, workload/time estimates) and preserves them across institutional data refreshes.
    \item The system supports reusable starting artefacts (e.g., template-based initiation) to enable peer-validated planning entry points.
\end{enumerate}

\paragraph{Delineation.}
Import from generic spreadsheet-like structures (e.g., CSV) is out of scope, since it does not affect the tool's own test scenario, in which curriculum data is loaded from a structured source rather than reconstructed from arbitrary user files.

\subsection{FEATURE-002 --- Multi-Level Roadmap \& Progress Dashboard}
\refstepcounter{featurenum}\label{feat:002}
\textbf{Importance score 19.} Associated themes: T1, T16, T18, T19, T57.

\paragraph{Description and user evidence.}
Participants plan at three simultaneous granularities, the current semester, the coming ones, and whole-programme progress \enquote{in percent or visually} (T1), and lack an always-current answer to \enquote{where am I right now?}, compensating by scrolling through transcripts (T16). Progress is conceptualised in credits or modules, \enquote{twelve of eighteen ECTS done} or \enquote{per module a bar} (T16, T57), and wanted prospectively: \enquote{if I add that, then in the future I would have 25\% fulfilled} (T18). A strict separation of \emph{planned} from \emph{completed} is described as essential, visually encoded, and motivating: \enquote{satisfying when it becomes less and less… I need a reward} (T19).

\paragraph{Requirements.}
\begin{enumerate}
    \item Users can navigate between Programme $\to$ Semester $\to$ Course levels.
    \item The Dashboard provides an always-current progress status including ECTS completed vs.\ required and percent completion at programme level.
    \item Progress can be presented module-based (e.g., per-module bar representation) in addition to aggregate progress.
    \item The system clearly distinguishes not planned vs.\ planned vs.\ completed states using consistent visual encodings (icons/colours/status labels) across table and graph views.
    \item Users can set a pacing target (e.g., target ECTS per semester) and the system displays a projected completion semester/range based on the current plan.
    \item Users can simulate the impact of adding/removing planned items and immediately observe updated progress (e.g., updated \%/ECTS and completion projection) and rewards.
\end{enumerate}

\paragraph{Delineation.}
The tool's minimum planning granularity is the semester, consistent with the overall scope (OOS2); course-level planning such as deadlines or scheduling collisions is not addressed.

\subsection{FEATURE-003 --- Curriculum Rules \& Compliance Engine}
\refstepcounter{featurenum}\label{feat:003}
\textbf{Importance score 18.} Associated themes: T13, T17, T30, T31, T33, T63.

\paragraph{Description and user evidence.}
Current artefacts do not operationalise curriculum rules, leaving the checking to the student: \enquote{Excel doesn't know your requirements… it doesn't detect errors… you have to double and triple check everything and hope it fits} (T13). Beyond category quotas, participants ask for dependency handling, \enquote{I'm missing… the dependencies between the subjects}, expecting warnings or blocks on invalid sequences (T31) and citing both hard pairs (\enquote{EP1 and EP2… they directly build on each other}, T31) and soft preferences (T30); trust rises when sequencing is institutionally grounded and marked as such (T63). They also want continuous fulfilment feedback, \enquote{you've already fulfilled five of seven requirements} (T17), and tie adoption to programme-specific tailoring (T33).

\paragraph{Requirements.}
\begin{enumerate}
    \item The system computes requirement fulfilment in real time at programme and category levels (e.g., \enquote{you've already fulfilled five of seven requirements}, \% completion), updating when a course is planned or marked completed.
    \item The system encodes hard rules, prerequisites, and sequencing constraints in real time and prevents scheduling that violates them (blocking).
    \item The system encodes soft/recommended prerequisites and issues a non-blocking warning that explains the recommendation (e.g., \enquote{not compulsory… but it would be good to do that one first}).
    \item The system identifies and lists what is missing (requirements/categories and actionable course options) and updates the missing list in real time.
    \item The compliance logic is shared across Graph and Table views, such that dependency structure and requirement status are consistent regardless of representation.
\end{enumerate}

\paragraph{Delineation.}
Institutional grounding of soft/recommended constraints is not available, since this information does not exist in the available data (OOS3).

\subsection{FEATURE-004 --- Graph-to-Table Planning Workflow}
\refstepcounter{featurenum}\label{feat:004}
\textbf{Importance score 15.} Associated themes: T20, T21, T44, T47, T61.

\paragraph{Description and user evidence.}
Participants describe a two-stage model, the graph for exploration and the table for committing, and are explicit that the activities are distinct: \enquote{graph used for exploration, not planning} (T20). Exploration is comprehensive and reassuring, \enquote{Show all… an immediate picture… satisfying} and \enquote{…I say plus, plus, plus… and then I've all the ones that interest me…}, valued for \enquote{less FOMO that I missed something again} (T20). Bridging the two calls for an intermediate holding space, \enquote{pool… with the courses… where I can move them around} (T21), with spatial consistency preserved, \enquote{like a desk… you just push things around} (T21), and committing for an explicit action with semester selection (T44). Table preferences favour familiarity, note-taking space (T47) and both manual and automatic ordering (T61).

\paragraph{Requirements.}
\begin{enumerate}
    \item Users can (a) browse/select courses in the Graph, (b) invoke Add to plan and select a target semester, and (c) see the course appear immediately in the Table under that semester with the same identity and status.
    \item The system provides a course pool/staging area that supports moving courses from exploration into semester planning.
    \item Drag-and-drop supports moving courses between pool and semesters and between semesters without data loss.
    \item Two-way synchronisation is implemented: changing a course's semester assignment in the Table updates the Graph state, and changes in Graph selection/planned state update the Table.
    \item The tool supports inline note-taking per course or per semester.
    \item The Table supports manual ordering and at least one automatic ordering mode (e.g., alphabetical or difficulty/priority), with the ability to revert to manual ordering after auto-sorting.
\end{enumerate}


\subsection{FEATURE-005 --- Feasibility Planning (Capacity, Workload Mix, Readiness)}
\refstepcounter{featurenum}\label{feat:005}
\textbf{Importance score 14.} Associated themes: T6, T35, T48, T49.

\paragraph{Description and user evidence.}
Course selection is often initially interest-driven (T6), but participants consistently frame planning \emph{quality} as bounded by feasibility: \enquote{at a glance… what is actually possible} (T6). Feasibility is anchored most concretely in ECTS load: participants target specific credit volumes (\enquote{so that I'm registered for around thirty}, T35) and describe iteratively reducing an over-ambitious plan (\enquote{31~ECTS… that's too much… brought it down… to about 21~ECTS}, T35). Beyond raw volume, workload \emph{mix} matters: participants balance project-based against exam-based courses and try to avoid exam clustering: \enquote{four exams in January… for me that would be too much} (T49). Finally, self-assessed readiness shapes selection: participants avoid courses that conflict with perceived weaknesses (\enquote{if you can't program well, then… maybe rather not take the programming-heavy subjects}, T48) and describe planning as an iterative, reflective process across semesters (T48).

\paragraph{Requirements.}
\begin{enumerate}
    \item Users can set semester capacity targets, including an ECTS target/range and optionally constraints such as maximum exams per assessment period.
    \item The system displays a semester feasibility summary that includes at minimum: total ECTS, distribution of assessment types, and flags when user-defined thresholds are exceeded.
    \item Users can record self-assessed readiness/strength dimensions (e.g., programming strength), and the system provides non-blocking guidance when planned courses are likely high-risk given that profile.
\end{enumerate}

\paragraph{Delineation.}
Requirement 3 was derived from the data and then scoped out of the implementation under the time and resource constraints of this thesis; it stands in this specification and is not withdrawn. Exam-load or exam-period collision detection is out of scope, since course-level planning and live institutional data are both excluded by the overall scope boundaries (OOS3, OOS2).

\subsection{FEATURE-006 --- Onboarding \& Guided Planning Transition}
\refstepcounter{featurenum}\label{feat:006}
\textbf{Importance score 12.} Associated themes: T5, T7, T59.

\paragraph{Description and user evidence.}
Planning needs shift over the study trajectory. Early semesters are typically followed by a prescribed curriculum: \enquote{the first semester was simply what was prescribed… the standard package with 29~ECTS} (T5). Autonomy then increases, particularly in the Master's programme (\enquote{In the Master's… only the first semester, and after that free choice}, T5), creating a recognisable transition point at which structured planning becomes necessary: \enquote{from around the third semester you realise: `Okay, now I need to start bringing some structure in'} (T7). Several participants report, in retrospect, that they would have benefited from adopting structured planning earlier (T7). For students new to university, the shift to self-managed planning is described as a high-friction adjustment: \enquote{with the whole change from technical secondary school to university… completely new to have to plan everything about it}, described by some as among the most difficult aspects of starting university (T59).

\paragraph{Requirements.}
\begin{enumerate}
    \item The product provides a guided first workflow: Show all $\to$ shortlist $\to$ add to semester plan, with contextual in-product tips shown the first time each step is encountered.
    \item Onboarding adapts guidance to user stage (e.g., early-semester \enquote{follow the standard package} vs.\ later-semester \enquote{build your own structure}), based on minimal inputs (e.g., current semester/programme stage).
\end{enumerate}


\subsection{FEATURE-007 --- Evidence-Based, Explainable Recommendation Support}
\refstepcounter{featurenum}\label{feat:007}
\textbf{Importance score 13.} Associated themes: T23, T24, T45, T52.

\paragraph{Description and user evidence.}
The default stance toward recommendations is sceptical, because the system cannot observe personal constraints: \enquote{I'm always very critical of recommendations, because the system can't know what my life is like… how many hours I work} (T23). Trust is described as emerging from inspectable evidence, a system like \enquote{a living, well-documented Wikipedia} (T23). Several explainable logics are endorsed: content- and dependency-based suggestions (T23), sequence guidance grounded in explicit relations (T45), and behavioural signals where available (T24). Peer and career-aligned suggestions are acceptable only as optional signals, \enquote{not as a main feature… for a reference point} (T23, T52), and autonomy is insisted on throughout: \enquote{I want to have the freedom} (T23).

\paragraph{Requirements.}
\begin{enumerate}
    \item The system generates recommendations using at minimum: (a) stated interest, (b) content similarity, and (c) dependency/sequence relations.
    \item Each recommendation is displayed with an explicit rationale label (e.g., matches interests, dependency/sequence relation, similar content, based on completed courses, internship lens).
    \item Each recommendation provides a minimally inspectable evidence view (e.g., matched tags, referenced dependency relation, summary/aggregate) and clearly indicates when evidence is available.
    \item Users can enable/disable recommendation families individually (content/interest, dependency/sequence, based-on-completed, peer-based, internship lens; behavioural signals only if available).
    \item Recommendations are non-blocking and do not prevent users from selecting any course; users can dismiss or hide recommendations without loss of core functionality.
    \item Users can apply a recommendation to shortlist or semester plan in one action while preserving graph/table context (viewport and selection preserved).
\end{enumerate}

\paragraph{Delineation.}
Recommendation signals based on scarcity or course-visitor data are excluded, since no live CMS/LMS data is available (OOS3). Peer-review-based recommendations are excluded, since no community feature is implemented (OOS4).

\subsection{FEATURE-008 --- Constraint-Aware Replanning \& Disruption Handling}
\refstepcounter{featurenum}\label{feat:008}
\textbf{Importance score 9.} Associated themes: T4, T46, T56.

\paragraph{Description and user evidence.}
Participants describe study plans as unstable, requiring frequent replanning in response to both personal and external disruption. Concurrent employment is a recurring destabilising constraint: \enquote{because I always worked alongside, it never worked out… I had to start shifting things} (T4), framed against a sustainability goal of managing studies \enquote{100\% and still be able to live a little} (T4). Externally, last-minute changes in course offerings are cited directly: \enquote{wasn't offered… shortly before the start of the semester}, with replacement courses being hard to find (T56). Participants also draw a pragmatic boundary around what a planning tool can reasonably be expected to solve, noting that frequent exam date/format changes are \enquote{not the system's problem} (T46).

\paragraph{Requirements.}
\begin{enumerate}
    \item Users can move a planned course to another semester, and the system immediately updates capacity, progress, and curriculum compliance/requirements status.
    \item Users can define a constraint profile (e.g., part-time status, max ECTS, pacing preferences) that drives feasibility summaries and warnings during replanning.
    \item When a course is unavailable (via data update or manual flag), the system indicates impacted requirements/categories, and provides alternative course suggestions where available.
    \item The system supports \enquote{quick replanning} interactions that preserve the existing plan structure (e.g., semester grouping and statuses) while enabling efficient substitution or rescheduling.
\end{enumerate}

\paragraph{Delineation.}
Automated detection of course unavailability is out of scope, since no live CMS/LMS connection is available (OOS3); unavailability must instead be reflected through manual flags or data updates.

\subsection{FEATURE-009 --- Overwhelm-Safe Browsing (Filters \& Abstraction)}
\refstepcounter{featurenum}\label{feat:009}
\textbf{Importance score 6.} Associated themes: T42, T50, T64.

\paragraph{Description and user evidence.}
Participants describe overload when confronted with the full catalogue or a dense graph: \enquote{way too much at the beginning} (T42), an \enquote{extreme number of courses thrown at you} that loses the overview (T64). The filter needs are stated concretely and cover status, category, ECTS range, obligation type and semester, with mandatory courses expected to stay visible by default (all T42). Beyond filtering they ask for abstraction, since \enquote{with graphs, it kind of explodes} (T50), proposing collapse and expand to control density (T42).

\paragraph{Requirements.}
\begin{enumerate}
    \item The Graph provides filters for at least: course status (not added / shortlisted / planned / completed), requirement type/category (e.g., mandatory, elective types), ECTS range, semester, and subject areas.
    \item Users can collapse/expand modules, groups, or the whole graph to reduce visible node count, while the system preserves context (which groups are collapsed and the current filter state).
    \item Filtering and collapsing do not delete user selections (shortlist/planned/completed) and do not unexpectedly reset saved layouts or positions.
\end{enumerate}


\subsection{FEATURE-010 --- Graph Interpretability (Purpose, Legend, Layout Modes)}
\refstepcounter{featurenum}\label{feat:010}
\textbf{Importance score 5.} Associated themes: T27, T32, T40, T41.

\paragraph{Description and user evidence.}
Participants report interpretability problems with the graph on first exposure: \enquote{at first glance I wouldn't know} its purpose (T41), and shorthand labels are \enquote{not intuitively understandable} (T40). Spatial layout compounds this: participants question the arrangement (\enquote{the layout confuses me… why such large padding}) and explicitly ask whether vertical position carries meaning (T27), a question the layout, as designed, does not answer visually. Participants express a reading-direction preference, \enquote{nicer if it goes left to right, because we read left to right}, and state the alternative to it in the same breath, \enquote{So either that, or I have the freedom to do it myself} (T32). Offering a free-arrangement option alongside the structured default is what lets both, a legible default layout and personal rearrangement, coexist rather than compete. The grouping and persistence side of free arrangement belongs to Feature~\ref{feat:013}, and its evidence is recorded there.

\paragraph{Requirements.}
\begin{enumerate}
    \item The Graph includes an always-available legend explaining symbols, colours, and status encodings, and stating what the current layout does or does not encode (including vertical and horizontal positioning).
    \item The tool offers at least two layout modes: (a) structured layout (left-to-right progression) and (b) free layout (user-arranged).
    \item On first use, the Graph View provides a short guided tooltip flow that enables users to correctly answer \enquote{what is this view for?}
    \item If the layout mode does not encode meaning (e.g., vertical or horizontal position), the UI explicitly communicates this to prevent misinterpretation.
\end{enumerate}


\subsection{FEATURE-011 --- Motivating Visual UI \& Emotional Load Reduction}
\refstepcounter{featurenum}\label{feat:011}
\textbf{Importance score 5.} Associated themes: T8, T54, T58.

\paragraph{Description and user evidence.}
Participants describe study planning as affectively demanding. Confronting remaining requirements can trigger distress: \enquote{very depressing, because only then did I see how much I still lack} (T8), including acute stress episodes during planning (T8). In parallel, participants describe visual clarity and polish as directly reducing perceived burden: a preference for low-text, icon-based interfaces (\enquote{with icons… not much text}; \enquote{too much text… off-putting}), self-identifying as visually oriented (\enquote{I'm… more of a visual person}), and linking an appealing interface to motivation to use it at all (all T54). The same low-clutter preference bears on completed items specifically: keeping every finished course fully rendered at full visual weight indefinitely would work against the low-text, uncluttered interface T54 calls for, without the record needing to be prominent once a course is done. Participants further indicate that the system should visibly take structuring work off their hands: \enquote{that the plan already takes quite a bit off your shoulders} (T58).

\paragraph{Requirements.}
\begin{enumerate}
    \item The UI supports a visual-first, low-text representation using icons and clearly distinguishable states with consistent colour meaning across views.
    \item Completed items can be faded and/or hidden while preserving an accessible record (i.e., completion remains visible on demand without deleting history).
    \item The system includes at least one positive reinforcement element tied to progress change (e.g., subtle completion animation or milestone indicator).
    \item Core planning actions are designed to reduce cognitive load (e.g., minimal required inputs, clear defaults, and immediate visual feedback after changes).
\end{enumerate}

\paragraph{Delineation.}
None beyond the overall scope boundaries (Section~\ref{sec:overall-scope}).

\subsection{FEATURE-012 --- Recommendation Presentation UX (Annotate, Don't Break Layout)}
\refstepcounter{featurenum}\label{feat:012}
\textbf{Importance score 4.} Associated themes: T43, T55.

\paragraph{Description and user evidence.}
Participants want recommendations to preserve the integrity of the exploration space rather than transform it: \enquote{marked… but not destroying the graph} (T43), explicitly contrasting removal-based filtering with lightweight marking, and naming \enquote{glowing borders} as a desired treatment (T43). At the same time, they are concerned about clutter, motivating a separate, controllable space for recommendation detail: \enquote{otherwise it's too much on one page}, framed as its own subsection for recommendations (T55). Participants also expect already-completed items to remain visible in context rather than being removed once a recommendation is shown (T43/T55).

\paragraph{Requirements.}
\begin{enumerate}
    \item Recommendations are displayed as annotations (e.g., overlays/badges/highlights such as borders) that do not trigger automatic graph re-layout or remove nodes by default.
    \item A dedicated Recommendation Panel/space exists that can be opened/closed without resetting the user's current graph viewport, selection, or filters.
    \item Users can accept a recommendation into shortlist or plan in one action, without losing current context (viewport remains stable; current selection preserved).
\end{enumerate}

\paragraph{Delineation.}
None beyond the overall scope boundaries (Section~\ref{sec:overall-scope}).

\subsection{FEATURE-013 --- Spatial Organisation \& Grouping Controls}
\refstepcounter{featurenum}\label{feat:013}
\textbf{Importance score 3.} Associated themes: T28, T29.

\paragraph{Description and user evidence.}
Participants describe free spatial arrangement as a cognitive externalisation strategy for expressing prioritisation, certainty, and personal relevance: \enquote{maybe I'd rather do it like that… because then I can place them how I want… draw in room to grow}, using vertical position to encode certainty (\enquote{up top… definitely… down here… I'm unsure}) and describing it as supporting their own mental grouping (all T28). In parallel, participants want grouping to be under their own control rather than system-imposed, asking directly how groups are formed and whether they can be changed (T29), and requesting the ability to freely create and name their own groupings: \enquote{the possibility to freely create groups… and name them how I want} (T29). They also raise the expectation that such groups should not be constrained by semester boundaries (T29).

\paragraph{Requirements.}
\begin{enumerate}
    \item Users can create named groups, move items into/out of groups, rename groups, and delete groups without deleting the underlying items.
    \item Spatial positions persist across sessions (no automatic reset that alters user-arranged layouts).
    \item Group assignments are visible and usable across contexts: group label is available in the Graph and is also surfaced in the Table (e.g., as a column/tag).
    \item Groups can include items spanning multiple semesters and do not break when an item's semester assignment changes.
\end{enumerate}

\paragraph{Delineation.}
Requirements 1, 3, and 4 were derived from the data and then scoped out of the implementation under the time and resource constraints of this thesis. Only free spatial positioning (Requirement 2) was built; the explicit, named, cross-view grouping mechanism was not. The requirements stand in this specification and are not withdrawn: what was bounded is what this thesis implemented, not what the evidence asked for.

\section{Traceability Matrix}
\label{app:traceability-matrix}
Each row names the Design chapter subsections that cite one feature. The themes each feature rests on and the number of requirements it carries are in the feature's own entry in Section~\ref{app:feature-specs} and are not repeated here. A section listed here means the feature is cited there, not that every one of its requirements is individually designed for: FEATURE-005 and FEATURE-013 each carry requirements that were not implemented, as their delineations record.

\small
\setlength{\tabcolsep}{4pt}
\begin{longtable}{@{}p{2.4cm}p{4.7cm}p{7.2cm}@{}}
\toprule
\textbf{Feature} & \textbf{Title} & \textbf{Design section(s)} \\
\midrule
\endfirsthead
\toprule
\textbf{Feature} & \textbf{Title} & \textbf{Design section(s)} \\
\midrule
\endhead
\bottomrule
\endfoot
FEATURE-001 & Unified Data \& Tool Integration & Curriculum Abstraction; Notes, Hours, and Grade Popover; Onboarding and First-Use Design \\
\addlinespace
FEATURE-002 & Multi-Level Roadmap \& Progress Dashboard & Course Lifecycle and State Machine; Module Backgrounds; Module Interaction; Planning Status Encoding; Dashboard Design; Text Content Across Hierarchy Levels \\
\addlinespace
FEATURE-003 & Curriculum Rules \& Compliance Engine & Compliance Engine Design; Curriculum Abstraction; Three-Panel Shell \\
\addlinespace
FEATURE-004 & Graph-to-Table Planning Workflow & Drag-and-Drop and Button Interaction Paths; Module Backgrounds; Module Interaction; Notes, Hours, and Grade Popover; Planning Workflow; Semester Lanes and Header Information; Semester Picker; Spatial Model of the Table View; Three-Panel Shell; Unified Interactive Canvas for Both Views \\
\addlinespace
FEATURE-005 & Feasibility Planning (Capacity, Workload Mix, Readiness) & Profile Design; Semester Lanes and Header Information \\
\addlinespace
FEATURE-006 & Onboarding \& Guided Planning Transition & Onboarding and First-Use Design; Planning Workflow \\
\addlinespace
FEATURE-007 & Evidence-Based, Explainable Recommendation Support & Accept Flow; Recommendation Families and Evidence; Recommendation Panel \\
\addlinespace
FEATURE-008 & Constraint-Aware Replanning \& Disruption Handling & Drag-and-Drop and Button Interaction Paths; Profile Design; Semester Picker \\
\addlinespace
FEATURE-009 & Overwhelm-Safe Browsing (Filters \& Abstraction) & Density Control and Multi-Dimensional Filtering \\
\addlinespace
FEATURE-010 & Graph Interpretability (Purpose, Legend, Layout Modes) & Left-to-Right Hierarchical Layout; Personal Customisability of Node Positions; Visual Legend Design \\
\addlinespace
FEATURE-011 & Motivating Visual UI \& Emotional Load Reduction & Context-Sensitive Action Buttons; Course Obligation Encoding; Planning Status Encoding; Dashboard Design; Subject Colour Hierarchy; Visual Design of Curriculum Elements \\
\addlinespace
FEATURE-012 & Recommendation Presentation UX (Annotate, Don't Break Layout) & Accept Flow; Recommendation Panel; Three-Panel Shell \\
\addlinespace
FEATURE-013 & Spatial Organisation \& Grouping Controls & Left-to-Right Hierarchical Layout; Module Interaction; Personal Customisability of Node Positions; Unified Interactive Canvas for Both Views \\
\end{longtable}
\normalsize